\chapter{团队与组织——谁在训练这些模型？}
\label{ch:org}
%% ============================================================

训练一个 frontier LLM 不仅是技术挑战，
更是一个复杂的组织工程。
它需要多个专业团队的紧密协作，
每个团队都面临独特的技术和管理挑战。
一个常见的误解是认为 LLM 训练的核心是算法研究——
事实上，工程执行占据了整个工作量的 80\% 以上。

\section{核心团队}

\subsection{数据团队：被低估的英雄}

数据团队是 LLM 训练中最被低估、但影响力最大的团队。
一个在业界广泛流传的经验法则是：
数据工程占据了整个训练工作量的 60--80\%。
模型架构可以从论文中复制，训练框架可以用开源工具，
但高质量的数据流水线必须从零构建。

数据团队的核心角色包括：

\begin{itemize}[noitemsep]
  \item \textbf{数据管道工程师}（Data Pipeline Engineers）：
        构建和维护从网页爬取到最终训练数据的全流程基础设施。
        这包括爬虫集群管理、数据清洗 pipeline（通常基于 Spark 或 Dask 构建）、
        去重系统（MinHash + LSH 在 PB 级数据上的高效实现）、
        以及数据版本管理（类似 Git LFS 但针对 TB 级数据集的版本控制）。
        这些工程师需要同时精通分布式系统和 NLP。
  \item \textbf{数据质量分析师}（Data Quality Analysts）：
        定义和维护数据过滤标准，什么样的文本是「高质量」的？
        这不是一个纯技术问题。
        perplexity 过低可能意味着模板化文本，过高可能意味着乱码。
        重复率、特殊字符比例、段落结构完整性：
        每一个过滤规则都需要在大量样本上验证，
        确保不会误杀有价值的数据。
  \item \textbf{语言学家}（Linguists）：
        对于多语言模型，语言学家的角色至关重要。
        分词器的设计需要考虑不同语言的形态学特征（中文的字/词粒度、
        日语的假名与汉字混合、阿拉伯语的从右到左书写和词根系统）。
        语言学家还负责评估不同语言的数据质量标准，
        适用于英语的 perplexity 过滤阈值不能直接套用到中文。
  \item \textbf{标注管理者}（Annotation Managers）：
        协调人类标注员团队，管理标注质量。
        在 SFT 和 RLHF 阶段，高质量的人类标注数据是稀缺资源。
        标注管理者需要设计标注指南、培训标注员、
        建立质量控制流程（inter-annotator agreement 监控、
        抽样审核、异常检测），
        并在成本和质量之间寻找平衡。
\end{itemize}

在一个主要 AI 实验室中，数据团队的规模通常在 10--30 人。
DeepSeek 在其技术报告~\cite{deepseekv3}中披露了
对数据配比的大量实验，网页、代码、学术论文、数学文本的最优比例
是通过在小模型上进行数百次 proxy 实验确定的。
这种工作量是纯算法研究无法覆盖的。

\begin{whybox}[为什么「数据质量 > 数据数量」？]
FineWeb~\cite{penedo2024fineweb} 的实验表明，
经过精细过滤的 1.3T tokens 数据集训练出的模型，
在多个基准测试上优于使用 4T tokens 低质量数据训练的模型。
直觉解释是：低质量数据不仅「浪费」了训练步数，
还可能引入噪声模式（如 HTML 标签、广告文本、机器翻译痕迹），
这些模式在推理时会以微妙的方式降低输出质量。
这就是为什么数据团队的工作如此关键。
\end{whybox}

\subsection{基础设施团队：让数千块 GPU 协同工作}

基础设施（Infra）团队负责 GPU 集群的管理与优化。
这个团队面临的核心挑战可以用一个数字概括：
在一个 16,384 块 GPU 的集群中，
假设每块 GPU 每天的故障率为 0.1\%，
那么平均每天会发生 $16384 \times 0.001 \approx 16$ 次 GPU 故障。
这意味着\textbf{平均每 90 分钟就有一块 GPU 出问题}。

Infra 团队的核心职责包括：

\begin{itemize}[noitemsep]
  \item \textbf{集群管理与任务调度}：
        使用 SLURM（传统 HPC 集群的标配）或 Kubernetes
        管理训练任务的调度。
        关键挑战是资源隔离，当多个训练任务共享集群时，
        如何确保网络带宽不互相干扰？
        如何在节点故障时自动重新调度而不中断其他任务？
  \item \textbf{故障检测与自动恢复}：
        这是 Infra 团队最核心的能力。
        自动检测 GPU 故障（CUDA error、ECC 错误、NVLink 断开）、
        网络故障（InfiniBand/RoCE 链路中断）、
        存储故障（NFS/Lustre 挂载点失效）。
        检测到故障后，系统需要自动隔离故障节点、
        将训练任务迁移到备用节点、
        从最近的 checkpoint 恢复训练，整个过程最好在 5--10 分钟内完成。
  \item \textbf{Checkpoint 管理}：
        对于 70B 参数的模型，每个完整训练 checkpoint 的大小约为
        $70 \times 10^9 \times 8 = 560\,\mathrm{GB}$
        （包含 BF16 参数、FP32 master weights 和 Adam 优化器状态 $m$/$v$，
        详见附录）。
        如果每 1000 步保存一次，保留最近 10 个 checkpoint，
        就需要 5.6\,TB 的高速存储空间。
        Infra 团队需要设计 checkpoint 的写入策略
        （异步写入，不阻塞训练）和存储架构
        （通常使用 Lustre 或 GPFS 这样的并行文件系统，
        提供数十 GB/s 的聚合带宽）。
  \item \textbf{网络监控与优化}：
        分布式训练的通信量巨大。
        在 tensor parallelism 中，
        每个 forward 和 backward pass 都需要 all-reduce 操作。
        Infra 团队需要监控网络拓扑中的热点、
        检测慢节点（stragglers）、
        优化 NCCL 的通信参数（如环形 vs.\ 树形 all-reduce、
        通信与计算的重叠度）。
\end{itemize}

\begin{corebox}[DeepSeek 的基础设施智慧]
DeepSeek~\cite{deepseekv3} 在 H800 集群（受出口管制限制，
NVLink 带宽低于 H100）上训练 V3 模型，
通过软件层面的创新弥补了硬件劣势：
DualPipe 通过更精细的计算-通信重叠减少了 pipeline bubble，
DeepEP 则为 MoE 的 expert parallelism
设计了专用的 all-to-all 通信内核。
他们实现了 “zero irrecoverable loss spikes”，
即整个训练过程中没有一次不可恢复的 loss 异常——
这需要极其成熟的容错基础设施。
\end{corebox}

\subsection{预训练研究团队：训练的「配方师」}

预训练研究团队负责回答一系列关键问题：
模型应该多大（多少层？多少注意力头？多大的 hidden dimension？）；
训练应该怎么做（学习率多少？batch size 多大？warmup 多长？数据怎么混合？）。

这个团队的工作方法是\textbf{先在小规模上验证，再推广到大规模}。
典型的工作流程：

\begin{enumerate}[noitemsep]
  \item 在 1B--7B 的 proxy model 上做超参数搜索，
        每组实验训练数小时到数天。
  \item 利用 Scaling Laws~\cite{kaplan2020scaling,hoffmann2022chinchilla}
        将小规模的最优配方推算到目标规模（如 70B 或 671B）。
  \item 在目标规模上做少量验证实验（训练数千步，
        观察 loss 曲线是否符合预期）。
  \item 确认配方后启动完整训练，通常持续数周到数月。
\end{enumerate}

「训练配方」（training recipe）是每个实验室的核心知识资产。
它包括但不限于：
\textbf{学习率调度}（cosine decay 的起点、终点和 warmup 步数）、
\textbf{batch size 调度}（从小 batch 逐渐增大，
还是全程使用固定 batch size？）、
\textbf{数据配比调度}（训练前期多用网页数据建立通用能力，
后期增加代码和数学数据提升推理能力）、
\textbf{长上下文扩展策略}（何时从 4K 切换到 128K，
RoPE 频率基数如何调整）。

LLaMA 3~\cite{dubey2024llama3} 的训练报告
详细披露了他们的训练配方：
学习率峰值 $1.5 \times 10^{-4}$，
cosine 衰减到峰值的 $1/30$，
warmup 8000 步，batch size 从 4M tokens 逐步增加到 16M tokens。
这些数字看似简单，但每一个都是大量实验的结果。

\subsection{后训练/对齐团队：定义「好」的含义}

后训练团队的工作是将一个「博学但不可控」的 base model
转变为一个「有用、无害、诚实」的 assistant。
这是整个训练流程中最依赖人类判断的环节。

核心职责包括：

\begin{itemize}[noitemsep]
  \item \textbf{RLHF/DPO 训练管道}：
        设计和实现偏好优化的训练流程。
        收集人类偏好数据（A 回答 vs.\ B 回答，哪个更好？），
        训练 reward model 或直接使用 DPO/GRPO 优化。
        关键挑战是偏好数据的一致性，
        不同标注员对「更好」的定义可能不一致。
  \item \textbf{Constitutional AI 设计}：
        受 Anthropic~\cite{bai2022constitutional} 启发，
        一些团队使用 AI 自身来生成和评估训练数据，
        基于一组预定义的「宪法原则」。
        这减少了对人类标注的依赖，
        但需要精心设计原则体系，
        什么是「有害的」？什么程度的帮助是「过度的」？
  \item \textbf{Red teaming 与安全评估}：
        组建专业的红队，系统性地测试模型的安全边界。
        红队会尝试越狱攻击（jailbreaking）、
        提示注入（prompt injection）、
        社会工程学手段（伪装成合法用户请求危险信息）。
        每一个发现的漏洞都需要通过额外的对齐训练来修复，
        然后再次测试，这是一个持续的迭代过程。
  \item \textbf{推理能力训练}：
        对于 reasoning model（如 DeepSeek-R1~\cite{deepseekr1}），
        后训练团队需要设计 RL 训练流程，
        让模型学会长链推理（chain-of-thought）、
        自我验证和错误修正。
        这需要精心设计的 reward signal，
        数学问题可以用答案正确性作为 reward，
        但开放性问题的 reward 设计要困难得多。
\end{itemize}

\begin{warnbox}[对齐的核心困难]
后训练团队面临的最根本挑战是：
\textbf{「有用、无害、诚实」这三个目标之间存在张力。}
过度追求无害会导致模型过度拒绝合理请求（over-refusal）；
过度追求有用会导致模型协助有害行为；
过度追求诚实会导致模型在不确定时拒绝给出有用的近似回答。
找到这三者的平衡点，本质上是一个价值观设计问题，
没有唯一正确的答案。
\end{warnbox}

\subsection{评估团队：度量你所关心的}

评估团队的工作是回答一个看似简单但极其困难的问题：
\textbf{模型到底有多好？}

核心职责包括：

\begin{itemize}[noitemsep]
  \item \textbf{Benchmark 套件管理}：
        维护覆盖多个能力维度的评估基准集合，
        数学推理（MATH、GSM8K）、代码生成（HumanEval、SWE-bench）、
        知识问答（MMLU）、指令遵循（IFEval）、
        长文本理解、多语言能力等。
        每个 benchmark 都需要持续维护，
        防止数据泄露（训练数据中包含了测试集）、
        更新过时的题目、处理评估格式变化。
  \item \textbf{内部评估管道}：
        构建自动化的评估基础设施。
        每次生成新的 checkpoint，
        系统需要自动在数十个 benchmark 上运行评估，
        生成对比报告。这本身就需要大量 GPU 资源，
        对 70B 模型在 MMLU 的 14K+ 题目上做推理
        可能需要数小时。
        高效的评估管道（如 lm-evaluation-harness 或 Lighteval）
        和评估集群的合理调度至关重要。
  \item \textbf{人类评估与 A/B 测试}：
        Benchmark 分数不能完全反映用户体验。
        评估团队还需要组织人类评估：
        让真实用户或专业评估者对比不同版本的模型输出，
        给出偏好判断。
        Chatbot Arena（LMSYS）的 Elo 评分系统
        是目前最受认可的人类评估方法。
  \item \textbf{评估的局限性}：
        评估团队最重要的认知是：
        \textbf{没有完美的评估}。
        模型可能在所有现有 benchmark 上表现优异，
        但在真实场景中出现令人意外的失败。
        评估团队需要持续设计新的评估方法，
        特别是针对 benchmark 无法捕捉的能力
        （如长期规划、常识推理、文化理解）。
\end{itemize}

\subsection{推理/部署团队：连接模型与用户}

推理/部署团队负责将训练好的模型变成用户可以使用的产品。
这个团队的工作直接决定了用户体验和运营成本。

核心职责包括：

\begin{itemize}[noitemsep]
  \item \textbf{服务架构设计}：
        API 网关 → 负载均衡器 → 推理集群 → 模型注册中心。
        需要处理突发流量（某个功能在社交媒体上火了，
        请求量可能在数分钟内增长 10 倍）、
        优雅降级（部分节点故障时不影响整体可用性）、
        自动扩缩容（根据负载动态调整推理服务器数量）。
  \item \textbf{延迟与成本优化}：
        用户期望的首 token 延迟（TTFT）通常在 500ms 以内，
        生成速度期望在 30--100 tokens/s。
        同时，推理成本直接影响定价和利润。
        部署团队需要在量化级别（FP8 vs.\ INT4）、
        批处理策略（continuous batching 的 max batch size）、
        缓存策略（prefix caching 的命中率优化）
        之间找到最优平衡。
  \item \textbf{SLA 管理}：
        商业 API 需要保证 99.9\% 以上的可用性
        （每月最多 43 分钟停机）。
        这需要多区域部署、灰度发布（新版本先在小流量上验证）、
        自动回滚机制。
  \item \textbf{滥用检测}：
        监控异常使用模式，大量请求尝试绕过安全限制、
        token farming（批量生成内容）、API key 泄露等。
        需要实时的异常检测系统和响应流程。
\end{itemize}

\section{组织模式：从实验室到公司}

观察最成功的 LLM 团队，
一个共同特点是\textbf{研究与工程的深度融合}——
写论文的人和写训练代码的人往往是同一批人。
但不同组织在规模、文化和战略上有显著差异。

\subsection{西方主要实验室深度解析}
\label{sec:org-labs}

\subsubsection{OpenAI：从非营利理想到商业巨擘}

\textbf{OpenAI}（2026 年初已超过 7,000 人）
经历了 AI 领域最戏剧化的组织演变。

\textbf{创立与理想主义时期（2015--2019）}。
2015 年 12 月，Sam Altman、Elon Musk、Greg Brockman 等人
以非营利组织的形式创立了 OpenAI，
承诺 10 亿美元的资助，
使命是「确保 AGI 惠及全人类」。
早期的 OpenAI 是一个纯粹的研究实验室，
发表了 GPT-1、GPT-2 等开创性论文，
但很快发现非营利模式无法支撑
训练 frontier 模型所需的天文数字般的计算投资。
Elon Musk 于 2018 年离开董事会
（后来他声称是因为对 OpenAI 商业化方向的不满），
这标志着理想主义时期的结束。

\textbf{商业化转型（2019--2023）}。
2019 年，OpenAI 创造了一种独特的
\textbf{capped-profit}（利润上限）结构——
投资者的回报被限制在 100 倍以内，
超出部分归属于非营利母体。
微软随即投入了 10 亿美元——
此后累计投资超过 130 亿美元。
GPT-3（2020 年）、ChatGPT（2022 年 11 月）
和 GPT-4（2023 年 3 月）的连续成功
使 OpenAI 从一个研究实验室
转变为全球最瞩目的 AI 公司。

\textbf{2023 年 11 月董事会危机}是 OpenAI 历史的转折点。
董事会突然解雇了 CEO Sam Altman，
原因至今没有完全公开，
但外界普遍认为与安全研究团队和商业化团队之间的
根本性路线分歧有关。
随后 95\% 以上的员工威胁集体辞职，
微软介入施压，Altman 在 5 天内复职。
这场危机的深层意义在于：
它暴露了 AI 实验室在
\textbf{安全使命}与\textbf{商业压力}之间的根本张力：
一个问题至今未得到解决，
只是暂时被增长的惯性所掩盖。

\textbf{Sam Altman 的领导风格}以
\textbf{极端的技术乐观主义}和
\textbf{超级推销员式的市场沟通}著称。
他在公开演讲中不断提高对 AGI 时间线的预期
（从「数十年」缩短到「可能几年内」），
同时积极推动 OpenAI 向全面商业化转型。
2025 年，OpenAI 在日本软银领投下
完成了 400 亿美元的融资，估值达到 3,000 亿美元。
Altman 于 2025 年 8 月任命前 Instagram 负责人
Fidji Simo 为 COO，并在 2026 年 2 月升任 CEO（应用），
自己则转为专注 AGI 研究和战略，
这一架构被解读为 Altman 对「研究突破」
比「商业运营」更感兴趣的信号。

\textbf{组织结构}高度垂直整合：
\textbf{Foundation Model 团队}（预训练与 scaling）、
\textbf{Post-Training 团队}（RLHF、推理训练）、
\textbf{Safety Systems 团队}（红队与对齐）、
\textbf{Applied AI/ChatGPT 团队}（产品工程）、
\textbf{Deep Research 团队}（长链推理与信息综合）。
从基础研究（GPT 系列、o 系列推理模型）
到消费产品（ChatGPT，月活超 4 亿）、
到企业 API、到硬件（自研芯片计划与 Broadcom 合作），
全部在一个屋檐下。

\textbf{人才流动}是 2025--2026 年 OpenAI 面临的显著挑战。
联合创始人 Ilya Sutskever 离开创立了
Safe Superintelligence Inc.（SSI），
多位安全团队核心成员出走，
一些人前往 Anthropic，另一些创立了新公司。
这种人才流失是否会影响 OpenAI 的长期竞争力，
取决于其继续吸引顶尖人才的能力，
而 3,000 亿美元的估值和 AGI 叙事
目前仍然是强大的人才磁铁。

\subsubsection{Anthropic：安全优先的 AI 实验室}

\textbf{Anthropic}（2025 年约 1,500 人，2026 年持续快速增长）
由前 OpenAI 研究副总裁 Dario Amodei (达里奥·阿莫德)
和他的姐姐 Daniela Amodei (达妮埃拉·阿莫德)
于 2021 年创立。

\textbf{创始 DNA 与安全使命}。
Dario Amodei 在 OpenAI 期间领导了 GPT-3 的训练，
是 scaling laws 的早期倡导者。
他离开 OpenAI 的原因是
对 OpenAI 在安全研究上的投入不足感到失望，
他认为随着模型能力的快速提升，
安全研究必须走在能力研究的\textbf{前面}而非后面。
Daniela Amodei 则负责 Anthropic 的商业和组织建设，
她的管理哲学是
「建立一个在追求安全的同时也能商业可持续的组织」，
安全不应该是一种慈善，而应该是一种竞争优势。

\textbf{Constitutional AI 作为核心哲学}~\cite{bai2022constitutional}。
与 OpenAI 依赖大量人类标注不同，
Anthropic 的对齐方法论以 Constitutional AI 为基石：
用一组明确的原则（「宪法」）
指导 AI 自身完成对齐训练。
这种方法的优势是\textbf{可扩展性}
（不需要线性增长的人类标注量）
和\textbf{一致性}
（原则是明确的，不受个别标注员偏好的影响）。
但它也依赖于一个哲学假设：
我们能够将「好的行为」编码为
一组有限的、可操作的原则。
这一假设在面对现实世界的道德复杂性时
是否成立，仍然是开放问题。

\textbf{Responsible Scaling Policy（RSP）}
是 Anthropic 独创的安全治理框架。
RSP 定义了一系列 \textbf{AI Safety Levels}（ASL），
类似于生物安全等级（BSL）：
ASL-1（无显著风险）到 ASL-4（潜在灾难性风险）。
模型在发布前必须通过对应级别的安全评估，
如果评估发现模型达到了更高的风险级别，
则必须实施更严格的安全措施
（如更强的访问控制、更多的红队测试）
才能继续部署。
RSP 的意义在于
它将安全承诺从模糊的「我们重视安全」
转化为\textbf{可验证的、有约束力的流程}。

\textbf{Claude 模型系列}的演进反映了 Anthropic 的技术路线：
2025 年推出了 Claude 3.7 Sonnet
（首个混合推理模型，同时支持快速回答和 extended thinking），
2026 年初发布了 Claude Opus 4 和 Claude Sonnet 4
（代号 Kryptonite 和 Borealis），
在编程和推理任务上取得了领先地位。
产品创新包括：
Claude Code（CLI 级别的自主编程 Agent）、
Computer Use（通过截图和键鼠操作桌面应用）、
以及 Model Context Protocol（MCP，Agent 工具交互的开放标准）。

Anthropic 在 2025 年完成了
Google、Salesforce 和 Amazon 领投的多轮融资，
估值超过 600 亿美元。

\subsubsection{Google DeepMind：规模与整合的挑战}

\textbf{Google DeepMind}（数千人，是规模最大的 AI 研究组织）
2023 年由 Google Brain 和 DeepMind 合并而成，
由 Demis Hassabis (德米斯·哈萨比斯) 领导。

\textbf{合并的动力与阵痛}。
Google Brain（2011 年成立，由 Jeff Dean 领导）
和 DeepMind（2014 年被 Google 收购，Hassabis 领导）
长期以来作为两个独立实体运营，
存在大量的重复研究和内部竞争。
两个团队各自发展了不同的技术栈
（Brain 使用 TensorFlow/JAX，DeepMind 使用自研框架），
不同的研究文化
（Brain 更偏工程，DeepMind 更偏理论），
甚至不同的发表策略。
2023 年的合并旨在消除这种内耗，
但整合过程并不平滑，
多位高级研究员在合并期间离职，
包括一些 Transformer 论文的原始作者。

\textbf{Ilya Sutskever 的离开}虽然发生在 OpenAI，
但其影响波及整个行业，
他的离开（以及 SSI 的创立）
激发了一波关于「安全 vs.\ 能力」路线辩论的浪潮，
也影响了 Google DeepMind 内部关于
安全研究资源分配的讨论。

\textbf{Gemini 系列}是 Google DeepMind 的旗舰产品。
Gemini 2.0 Flash（2024 年底）引入了原生工具使用和代码执行，
Gemini 2.5 Pro（2025 年初）成为首个
在 LMSYS Chatbot Arena 登顶的 Google 模型，
而 Gemini 3（预计 2026 年中发布）
将进一步推进原生多模态和长上下文能力。

Google DeepMind 的\textbf{独特优势}是
\textbf{端到端的基础设施控制}：
从芯片设计（TPU v5、v6）到网络架构到训练框架（JAX/Flax）
到云服务（Google Cloud）再到消费产品
（Gemini App、Google Search 的 AI Overviews），
全栈自研。
这种垂直整合理论上提供了最大的优化空间，
但也带来了「大公司病」的风险，
决策链条长、产品发布慢、
多个产品团队（Gemini App、Workspace AI、Cloud AI）
有时会出现内部竞争和资源争夺。

\subsubsection{Meta FAIR：开源作为竞争策略}

\textbf{Meta FAIR 及 GenAI 团队}（合计 1,000+ 人）。
Mark Zuckerberg 做出了 AI 行业最大胆的战略选择：
将 LLaMA~\cite{touvron2023llama,dubey2024llama3} 系列开源。

\textbf{开源的战略逻辑}不是慈善，而是一种精心计算的竞争策略。
当你的模型成为行业标准，围绕它的生态系统（工具、微调模型、应用）
就成为你的护城河。Meta 不靠卖 API 赚钱，
而是通过 AI 能力提升其广告和社交产品的收入：
如果 LLaMA 成为大多数企业的 AI 基础设施，
Meta 就间接控制了 AI 应用层的技术标准。
2024 年 Zuckerberg 发表的公开信
明确阐述了这一逻辑：
「开源 AI 对 Meta 有利，因为它推动了生态系统的增长，
而生态系统的增长有利于我们的产品。」

\textbf{LLaMA 4 系列}（2025 年 4 月发布，
包括 Scout 109B/17B-active 和 Maverick 400B/17B-active）
采用了 MoE 架构，是 Meta 开源的首个 MoE 模型。
Yann LeCun (杨立昆) 仍担任 Chief AI Scientist，
负责 FAIR 的长期研究方向。
LeCun 长期以来对自回归语言模型持批评态度，
倡导一种基于\textbf{世界模型}（world model）
和\textbf{联合嵌入预测架构}（JEPA）的替代路线，
他认为当前的 LLM 只是「鹦鹉学舌」，
缺乏真正的理解和规划能力。
这种内部的思想张力
（LeCun 的长期愿景 vs.\ GenAI 团队的产品压力）
是 Meta AI 组织动力学的一个有趣特征。

\textbf{Research SuperCluster}（RSC）
是 Meta 的 AI 训练基础设施，
截至 2026 年初拥有超过 60 万块 GPU，
是全球最大的 AI 训练集群之一。
Meta 在 2025 年宣布计划在 2026 年将
GPU 总量扩展至 130 万块，
资本支出预算高达 600--650 亿美元。
这种规模的投资使得 Meta 在算力上
与 Google 和微软处于同一量级。

\subsubsection{xAI：Elon Musk 的 AI 野心}

\textbf{xAI}（2023 年 7 月创立）
由 Elon Musk 在离开 OpenAI 董事会后创立，
创始成员包括来自 DeepMind、Google Brain
和多伦多大学的研究者。

\textbf{Musk 的创立动机}是多层面的。
表面上，他声称是为了对抗 OpenAI 和 Google 的
「AI 安全不足」和「过度政治正确」；
深层上，AI 能力对 Musk 的商业帝国
（Tesla 自动驾驶、SpaceX 任务规划、X/Twitter 内容推荐）
具有战略意义。
xAI 从一开始就定位为
一个同时服务研究和产品的组织。

\textbf{Memphis 超级集群}是 xAI 的硬件基础。
2024 年，xAI 在田纳西州 Memphis 建设了
由 \textbf{100,000 块 H100 GPU} 组成的超级集群，
建设速度创下了记录（从动工到投入使用仅约 19 天的内部装配）。
这个集群的规模使 xAI 在算力上
直接与 OpenAI 和 Google 竞争。
2025 年底，Musk 宣布计划将其扩展至
\textbf{200,000 块 GPU}——
进一步巩固 xAI 的算力优势。

\textbf{Grok 系列}是 xAI 的模型产品线。
Grok 1（2023 年底）作为 X 平台的内置 AI 助手发布，
Grok 2（2024 年）在多个基准上达到 GPT-4 水平，
Grok 3（2025 年初）引入了推理能力。
Grok 的独特卖点是与 X（前 Twitter）的实时数据集成，
模型可以访问最新的推文和新闻，
提供其他模型无法获取的实时信息。
但这也带来了独特的安全挑战，
实时数据中可能包含大量错误信息和有害内容。

\subsubsection{Mistral AI：欧洲的 AI 旗帜}

\textbf{Mistral AI}（约 600--700 人，2025 年）
是欧洲最重要的 AI 公司，
由 Arthur Mensch（前 Google DeepMind）、
Guillaume Lample（前 Meta）和
Timoth\'{e}e Lacroix（前 Meta）于 2023 年创立。

Mistral 的创立故事本身就是一个产业叙事：
三位创始人都来自全球最顶尖的 AI 实验室，
但选择回到巴黎创业，
部分原因是欧洲提供了不同于硅谷的
监管环境和人才池，
部分原因是他们相信
\textbf{效率优先}的路线
可以在不追求最大规模的情况下
产出高质量的模型。

\textbf{Mixtral}~\cite{jiang2024mixtral} 是 Mistral 的标志性成果：
一个在 2024 年初发布的 MoE 模型，
证明了一个精干的团队可以在 MoE 架构上
做出与 GPT-3.5 级别竞争的工作。
随后的 Mistral Large 2 和 Pixtral（多模态）
进一步巩固了 Mistral 的技术地位。

Mistral 在欧洲市场的\textbf{独特竞争优势}是
\textbf{EU AI Act 合规}。
作为欧洲公司，Mistral 天然更理解
也更容易满足欧洲的 AI 监管要求，
这对需要合规部署的欧洲企业客户
是一个重要的差异化卖点。
Mistral 为这些客户提供了
私有部署方案和定制微调服务，
在企业 AI 市场建立了稳固的根据地。

但 Mistral 也面临挑战：
作为一家从研究导向转向商业化的公司，
它在 2025 年从 50--80 人快速扩张到 600--700 人，
这种增长速度对组织文化和执行效率都是考验。
在算力和资金上，Mistral 与美国和中国的顶级玩家
仍有一到两个数量级的差距，
这迫使它必须在效率上做到极致
才能保持竞争力。

\subsection{中国 AI 实验室生态}
\label{sec:chinese-labs}

中国的 AI 实验室在 2024--2026 年经历了
从「追赶者」到「并行创新者」的关键转变。
DeepSeek-R1 的成功标志着
中国团队在某些维度上
已经能够与美国顶尖实验室正面竞争。
但中国 AI 生态的复杂性远超单一机构的成就，
它是一个由数十个互相竞争、各具特色的实验室
组成的\textbf{多极格局}。

%% --- TikZ: Global AI Lab Competitive Landscape ---
\begin{figure}[htbp]
\centering
\begin{tikzpicture}[
  >={Stealth[length=2.5pt,width=2.5pt]},
  every node/.style={font=\small},
  regionbox/.style={rectangle, rounded corners=5pt, draw=#1!40,
    fill=#1!4, inner sep=6pt, minimum width=48mm, line width=0.6pt},
  regiontitle/.style={font=\small\bfseries, #1},
  labfrontier/.style={font=\scriptsize, inner sep=0pt, #1},
  labother/.style={font=\tiny, figGray, inner sep=0pt},
  metricstyle/.style={font=\tiny, figGray!70},
]

% ======= US Column =======
\node[regionbox=figBlue, minimum height=68mm, anchor=north] (usbox) at (-5.2, 0) {};
\node[regiontitle=figBlue, anchor=north] at (usbox.north) [yshift=-3mm] {美国};
\draw[figBlue!30, line width=0.4pt] (-6.9, -0.9) -- (-3.5, -0.9);

% Frontier labs
\node[labfrontier=figBlue, anchor=north west] at (-6.8, -1.2)
  {\textbf{OpenAI}};
\node[metricstyle, anchor=north west] at (-6.8, -1.55)
  {7,000+ 人 · \$13B 融资};
\node[labfrontier=figBlue, anchor=north west] at (-6.8, -2.1)
  {\textbf{Anthropic}};
\node[metricstyle, anchor=north west] at (-6.8, -2.45)
  {1,500+ 人 · \$8B 融资};
\node[labfrontier=figBlue, anchor=north west] at (-6.8, -3.0)
  {\textbf{Google DeepMind}};
\node[metricstyle, anchor=north west] at (-6.8, -3.35)
  {Gemini 系列 · TPU v5};
\node[labfrontier=figBlue, anchor=north west] at (-6.8, -3.9)
  {\textbf{Meta FAIR}};
\node[metricstyle, anchor=north west] at (-6.8, -4.25)
  {1,000+ 人 · Llama 开源};

% Divider
\draw[figBlue!20, line width=0.3pt, dashed] (-6.8, -4.7) -- (-3.6, -4.7);
\node[font=\tiny, figBlue!50, anchor=west] at (-6.8, -4.9) {其他};

\node[labother, anchor=north west] at (-6.8, -5.2) {xAI · 100K H100};
\node[labother, anchor=north west] at (-6.8, -5.6) {AI2 · OLMo 开源};
\node[labother, anchor=north west] at (-6.8, -6.0) {Cohere · 企业级};

% ======= China Column =======
\node[regionbox=figRed, minimum height=68mm, anchor=north] (cnbox) at (0, 0) {};
\node[regiontitle=figRed, anchor=north] at (cnbox.north) [yshift=-3mm] {中国};
\draw[figRed!30, line width=0.4pt] (-1.7, -0.9) -- (1.7, -0.9);

% Frontier labs
\node[labfrontier=figRed, anchor=north west] at (-1.6, -1.2)
  {\textbf{DeepSeek}};
\node[metricstyle, anchor=north west] at (-1.6, -1.55)
  {150--200 人 · 效率优先};
\node[labfrontier=figRed, anchor=north west] at (-1.6, -2.1)
  {\textbf{月之暗面}};
\node[metricstyle, anchor=north west] at (-1.6, -2.45)
  {Kimi · 长上下文};
\node[labfrontier=figRed, anchor=north west] at (-1.6, -3.0)
  {\textbf{智谱 AI}};
\node[metricstyle, anchor=north west] at (-1.6, -3.35)
  {ChatGLM · 全栈};
\node[labfrontier=figRed, anchor=north west] at (-1.6, -3.9)
  {\textbf{字节豆包}};
\node[metricstyle, anchor=north west] at (-1.6, -4.25)
  {Seed 系列 · 端侧};

% Divider
\draw[figRed!20, line width=0.3pt, dashed] (-1.6, -4.7) -- (1.6, -4.7);
\node[font=\tiny, figRed!50, anchor=west] at (-1.6, -4.9) {其他};

\node[labother, anchor=north west] at (-1.6, -5.2) {01.AI · Yi 系列};
\node[labother, anchor=north west] at (-1.6, -5.6) {阶跃星辰 · 视频生成};
\node[labother, anchor=north west] at (-1.6, -6.0) {MiniMax · 消费产品};
\node[labother, anchor=north west] at (-1.6, -6.4) {Qwen · 开源生态};

% ======= EU Column =======
\node[regionbox=figGreen, minimum height=68mm, anchor=north] (eubox) at (5.2, 0) {};
\node[regiontitle=figGreen, anchor=north] at (eubox.north) [yshift=-3mm] {欧洲};
\draw[figGreen!30, line width=0.4pt] (3.5, -0.9) -- (6.9, -0.9);

% Frontier lab
\node[labfrontier=figGreen, anchor=north west] at (3.6, -1.2)
  {\textbf{Mistral}};
\node[metricstyle, anchor=north west] at (3.6, -1.55)
  {600--700 人 · \$2B+ 估值};
\node[metricstyle, anchor=north west] at (3.6, -1.95)
  {开源 + 商业双轨};

% Divider
\draw[figGreen!20, line width=0.3pt, dashed] (3.6, -2.5) -- (6.8, -2.5);
\node[font=\tiny, figGreen!50, anchor=west] at (3.6, -2.7) {生态};

\node[labother, anchor=north west] at (3.6, -3.0) {Aleph Alpha (德国)};
\node[labother, anchor=north west] at (3.6, -3.4) {Stability AI (英国)};

\end{tikzpicture}
\caption{全球 AI 实验室竞争格局（2026 年初）。
美国仍在规模和资金上占据优势，
中国在效率创新和开源贡献上形成了独特的竞争力，
欧洲以 Mistral 为代表建立了第三极。}
\label{fig:ai-lab-landscape}
\end{figure}

\subsubsection{DeepSeek：量化基金的 AGI 野心}

\textbf{DeepSeek}（约 150--200 名核心研究员和工程师）
背靠量化对冲基金\textbf{幻方量化}（High-Flyer Capital Management），
是「小团队大影响」的典范。

\textbf{幻方量化的 DNA}深刻塑造了 DeepSeek 的组织文化。
幻方量化是中国最成功的量化对冲基金之一，
管理资产规模曾超过 600 亿人民币。
量化基金的核心竞争力是
用极少的人、极高的效率、极精密的系统
在金融市场中寻找微小的套利机会。
这种文化——对效率的极致追求、
对系统工程的高度重视、
对冗余人员的零容忍——
被完整地移植到了 DeepSeek 的 AI 研究中。

\textbf{梁文锋 (Liang Wenfeng)}
是幻方量化和 DeepSeek 的创始人。
他是浙江大学电子信息工程本科毕业，
之后进入量化交易行业。
与大多数 AI 实验室的创始人不同
（通常来自学术界顶尖 AI 实验室），
梁文锋的背景是\textbf{工程和金融}。
这种背景赋予了他两个独特优势：
第一，\textbf{系统工程思维}——
量化交易系统对延迟、吞吐量和可靠性的要求
与 AI 训练系统高度相似，
幻方积累的分布式系统经验
直接加速了 DeepSeek 的基础设施建设；
第二，\textbf{资金独立性}——
幻方的自有资金使 DeepSeek 不需要依赖外部风投，
因此不受短期商业化压力的驱动，
可以专注于长期技术突破。

\textbf{为什么一家量化基金要建 AGI？}
梁文锋在接受采访时给出的回答是：
「我们在量化交易中使用了大量的机器学习技术，
逐渐意识到通用智能才是最终的工具。
如果你能造出足够聪明的 AI，
量化交易只是它能做的无数件事之一。」
这个回答透露了一种\textbf{第一性原理}的思维方式，
不是为了解决具体问题而做 AI，
而是因为相信 AGI 是最终的杠杆。

\textbf{DeepSeek 的技术轨迹}异常陡峭：
\begin{itemize}[noitemsep]
  \item \textbf{DeepSeek-V2}（2024 年 5 月）：
        引入 MLA（Multi-head Latent Attention）架构，
        将 KV cache 压缩了 10 倍以上。
  \item \textbf{DeepSeek-V3}~\cite{deepseekv3}（2024 年 12 月）：
        671B 参数 MoE 模型，
        训练成本仅 \$5.5M（GPU 租赁），
        达到 GPT-4 级别性能。
  \item \textbf{DeepSeek-R1}~\cite{deepseekr1}（2025 年 1 月）：
        通过纯 RL 训练实现推理能力涌现，
        被称为「中国的 Sputnik 时刻」。
        2025 年 9 月登上 \textit{Nature} 封面。
  \item \textbf{V4 开发中}（2026 年）：
        据报道聚焦于原生多模态、
        LTM（Long-Term Memory）突破、
        以及国产芯片（华为 Ascend）适配优化。
\end{itemize}

DeepSeek 用不到 Meta FAIR 十分之一的人员，
达到了 GPT-4 级别的性能。
关键因素包括：
极致的工程效率（MLA、DeepEP 等创新）、
扁平的组织结构（减少管理层级，加速决策）、
以及量化基金背景带来的对效率的极度重视。
团队结构高度集成，
没有独立的「数据团队」「Infra 团队」「研究团队」的明确边界，
研究员同时负责算法设计和底层系统优化。

\subsubsection{月之暗面 (Moonshot AI)：长上下文的先行者}

\textbf{月之暗面}由\textbf{杨植麟 (Yang Zhilin)} 于 2023 年创立。
杨植麟本科就读于清华大学，
博士在 CMU（师从 Ruslan Salakhutdinov 和 William Cohen），
之后在 Google Brain 和 Meta AI 工作。
他最广为人知的学术贡献是
\textbf{Transformer-XL}~\cite{dai2019transformerxl}
和 \textbf{XLNet}~\cite{yang2019xlnet}——
两篇在长序列建模领域具有开创性的论文。

这种学术背景直接影响了月之暗面的技术路线：
从一开始，\textbf{Kimi} 就以
\textbf{超长上下文}能力作为核心差异点。
2024 年初，Kimi 率先支持 200K token 上下文，
远超同期竞品（大多为 32K--128K），
引发了中国 AI 市场的「上下文长度竞赛」。

Kimi 的产品策略是
\textbf{技术差异化 + 知识工作者定位}。
与豆包、文心一言等面向大众市场的产品不同，
Kimi 定位于专业用户，
学者需要阅读和理解长篇论文、
律师需要分析复杂合同、
开发者需要理解大型代码库。
这种定位使得 Kimi 在
一个有消费意愿的细分市场中建立了品牌认知，
虽然总用户量不如字节或百度的产品，
但用户黏性和付费率更高。

2025 年底，月之暗面推出了 \textbf{Kimi K2.5}——
支持 Agent 能力，
在 AIME 2025 上达到 96.1\%，
证明长上下文优势可以扩展到推理能力。

\subsubsection{智谱AI (Zhipu AI)：学术到产业的桥梁}

\textbf{智谱AI}由清华大学教授\textbf{唐杰 (Tang Jie)} 联合创立，
是中国学术界与 AI 产业之间
最成功的技术转化案例之一。

唐杰是清华大学计算机系教授，
在学术图谱（AMiner）和图神经网络领域有深厚积累。
智谱AI 与清华大学知识工程实验室（KEG）深度绑定，
这种学术根基赋予了它两个独特优势：
第一，\textbf{人才供给}——
清华 CS 每年的毕业生成为智谱的人才蓄水池；
第二，\textbf{技术传统}——
GLM（General Language Model）架构
是在清华内部发展起来的，
与 GPT 的纯因果语言模型不同，
GLM 使用了\textbf{自编码 + 自回归}的混合训练目标，
在某些理解任务上具有理论优势。

\textbf{ChatGLM 系列}的演进：
ChatGLM-1（2023 年初，中国最早的开源对话模型之一），
ChatGLM-2、ChatGLM-3（持续迭代），
GLM-4（2024 年），
GLM-4.7（2026 年初，融合多模态和长上下文）。
GLM 系列的一个特点是
\textbf{持续坚持 GLM 架构}而非切换到纯 GPT 架构，
这在技术上是一种冒险
（因为 GPT 架构有更成熟的训练生态），
但也形成了技术差异化。

智谱AI 在 2025 年获得了超过 40 亿人民币的融资，
投后估值超过 300 亿人民币，
是中国估值最高的 AI 创业公司之一。

\subsubsection{零一万物 (01.AI)：Kai-Fu Lee 的「大而快」战略}

\textbf{零一万物}由\textbf{李开复 (Kai-Fu Lee)} 于 2023 年创立。
李开复是 AI 领域最知名的华人面孔之一，
曾任 Google 中国区总裁，
后创办创新工场专注 AI 投资，
著有《AI 超级大国》等畅销书。

\textbf{Yi 系列}是零一万物的模型产品线。
Yi-1（2023 年底）是中国最早的开源基座模型之一，
Yi-1.5（2024 年）引入了 MoE 架构，
Yi-Lightning（2024 年底）定位为高速推理模型。
零一万物采用了\textbf{开放权重}策略，
模型权重开源，但训练数据和训练代码不公开。

李开复的战略理念是
「\textbf{追赶最前沿的模型}不如\textbf{做最适合中国市场的应用}」，
他很早就意识到，
在基座模型层面与 OpenAI 和 Google 正面竞争
可能不是最优策略，
而是应该在应用层找到差异化定位。
这种务实的思维使得零一万物
在商业化方面比一些纯研究导向的实验室更快，
但也使它在技术影响力上
不如 DeepSeek 或月之暗面那样引人注目。

2025 年下半年，零一万物经历了战略调整：
从「做中国最好的 base model」
转向「做中国最好的 AI 应用平台」，
模型层投入相对减少，应用层投入加大。

\subsubsection{阶跃星辰 (StepFun)：从语言到视频}

\textbf{阶跃星辰}由\textbf{姜大昕 (Jiang Daxin)} 创立。
姜大昕曾任微软亚洲研究院（MSRA）副院长，
在自然语言处理和信息检索领域有数十年经验。
他从 MSRA 离职创业时带走了一批核心研究人员，
构成了阶跃星辰的技术骨干。

阶跃星辰的独特之处在于
它是中国 AI 创业公司中
\textbf{最早大规模投入视频生成}的之一。
Step 系列模型从语言模型起步
（Step-1V、Step-2，追赶 frontier 水平），
但 2025 年开始将重心转向\textbf{视频生成和理解}——
Step Video 在视频质量和一致性上
达到了与 Sora 和 Kling 竞争的水平。

\textbf{Step 3.5 Flash}（2026 年初）
在 AIME 2025 上达到 97.3\%，
展示了阶跃星辰在推理能力上的快速追赶。
但公司的长期赌注是
\textbf{多模态}——特别是视频：
将成为下一代 AI 应用的核心交互方式。

\subsubsection{MiniMax：消费产品的先驱}

\textbf{MiniMax} 由\textbf{闫俊杰 (Yan Junjie)} 创立。
闫俊杰曾是商汤科技（SenseTime）的技术副总裁，
在计算机视觉和多模态领域有深厚积累。

MiniMax 的策略与大多数中国 AI 实验室截然不同：
它从一开始就\textbf{以消费产品为导向}——
而非以基座模型为导向。
其旗舰产品包括：
\textbf{Talkie}（海外版 AI 社交/角色扮演 App，
在美国市场获得了数百万下载量）
和 \textbf{海螺AI/Hailuo}（视频生成产品）。

MiniMax 的逻辑是：
基座模型之间的差距在快速缩小，
长期来看，模型层可能商品化，
真正的壁垒在于\textbf{产品体验和用户关系}。
Talkie 的成功部分验证了这一假设：
它使用的并非最强的基座模型，
但通过出色的产品设计
（角色一致性、情感交互、社区运营）
建立了用户黏性。

\subsubsection{字节跳动豆包 (ByteDance Doubao)：分发为王}

\textbf{字节跳动}凭借 TikTok/抖音的全球分发能力，
在 AI 产品层拥有独特的结构性优势。

\textbf{豆包}是字节跳动的 AI 助手产品，
得益于抖音生态的流量灌注，
其用户增长速度远超独立 AI 创业公司。
但豆包的技术核心，\textbf{Seed 团队}：
在基座模型层面一直保持着相对低调的姿态。

2026 年初，Seed 团队通过积极的人才吸纳
构建多模态 AI 的全面能力。
特别值得注意的是：
2026 年 3 月，\textbf{Qwen 大模型后训练负责人郁博文}
从阿里巴巴离职加入 Seed 团队，
负责视觉模型和多模态交互方向的后训练工作。
Seed 团队的 \textbf{Doubao Seed 2.0 Pro}
在 AIME 2025 上达到 98.3\%，
这是中国团队在数学推理基准上的最高分，
证明了字节在基座模型层面的实力
远超外界的普遍认知。

字节跳动的核心优势不在于
单一模型的基准分数最高，
而在于\textbf{从模型到用户的通路最短}：
一个新模型可以在数天内
通过抖音、今日头条、飞书等产品
触达数亿用户。
这种分发能力使得字节在
「模型 $\to$ 产品 $\to$ 用户反馈 $\to$ 模型改进」
的飞轮上具有结构性优势。

\subsubsection{阿里巴巴 Qwen：开源多语言的领军者}

\textbf{阿里巴巴通义}团队开发的 \textbf{Qwen 系列}~\cite{qwen2025qwen3}
是中国开源模型中国际影响力最大的系列之一。

Qwen 的核心优势是\textbf{多语言能力}——
特别是在中文和英文之间的平衡。
多数模型要么以英文为主（如 LLaMA），
要么以中文为主（如早期的 ChatGLM）。
Qwen 通过精心设计的
多语言训练数据配比和分词器优化，
在中英双语上同时达到了 frontier 水平。
Qwen3（2025 年）采用 Apache 2.0 许可证，
进一步降低了商用门槛。
\textbf{Qwen3-Omni}（2025 年底）
是中国首个原生多模态大模型之一，
支持文本、图像、语音的统一理解和生成。

但 Qwen 团队在 2026 年面临的挑战同样严峻。
核心人才的流失（后训练负责人郁博文等离职）、
阿里巴巴内部从研究向产品的战略转移、
以及来自 DeepSeek 和字节的激烈竞争，
都在考验这支团队的凝聚力和执行力。
阿里云的 Token Hub 战略
（将 Qwen 模型的推理能力商品化，
以极低价格提供 API 服务）
可能削弱了基础研究投入的激励：
当你的商业模式是「卖 token」而非「做最好的模型」时，
研究团队的使命感会面临微妙的挑战。

\subsubsection{中国 AI 生态的结构性特征}

观察中国 AI 实验室的整体格局，
有几个结构性特征值得注意：

\textbf{资金来源的多样性}。
与美国 AI 创业公司主要依赖 VC 融资不同——
中国 AI 实验室的资金来源更加多元：
DeepSeek 背靠量化基金的自有资金，
智谱AI 有清华大学的资源支持，
月之暗面获得了红杉中国等顶级 VC 的投资，
而字节和阿里的 AI 投入来自各自的核心业务利润。
这种多元性意味着中国 AI 的发展
不会因为单一资金来源的枯竭而受阻。

\textbf{「内卷」式竞争的利弊}。
中国 AI 市场的竞争烈度远超美国，
十几家主要实验室在同一个赛道上
争夺相似的人才、数据和算力资源。
这种竞争推动了快速进步
（中国模型从落后 GPT-4 两年
到与之平齐只用了约 18 个月），
但也带来了资源浪费
（大量重复的训练实验、
价格战导致的利润率压缩）。
一些行业观察者认为，
到 2027 年可能会出现
显著的\textbf{整合}——
资金不足的二三线实验室
将被收购或关闭，
最终形成 3--5 个主要玩家的格局。

\textbf{应用创新的差异化路径}。
与美国更注重基础模型能力不同，
中国 AI 生态的一个显著特点是
\textbf{应用层的快速创新}。
从 AI 社交（Talkie/星野）
到 AI 教育（Kimi 的论文阅读）
到 AI 短视频（豆包+抖音）
到 AI 电商（通义千问+淘宝），
中国团队在将 AI 嵌入消费场景方面
展现了令人瞩目的速度和创造力。
这种「应用驱动」的创新模式
与美国「模型驱动」的创新模式
形成了有趣的互补，
双方正在不同的维度上推进 AI 的边界。

\begin{keybox}[「小团队，大影响」vs.「大平台，快分发」]
中国 AI 生态中最有趣的对比是
DeepSeek 式的\textbf{精英小团队模式}
与字节跳动式的\textbf{大平台分发模式}。

DeepSeek（150--200 人）的优势在于：
极高的人才密度、极快的决策速度、
极低的管理开销、
以及不受商业化压力驱动的研究自由度。
它用不到 Meta FAIR 十分之一的人员
产出了 V3/R1。

字节跳动 Seed 团队的优势在于：
海量的用户数据飞轮、
全球级的产品分发能力、
近乎无限的计算预算、
以及吸引顶尖人才的薪资竞争力。

两种模式各有所长，
DeepSeek 更适合\textbf{技术突破}——
字节更适合\textbf{规模化落地}。
但长期来看，两者可能会收敛：
DeepSeek 在 V4 之后也需要面对产品化问题，
字节也在持续提升基础研究的深度。
\end{keybox}

\section{Jensen Huang 的「五层」愿景}
\label{sec:jensen-five-layers}

NVIDIA CEO 黄仁勋 (Jensen Huang) 在多次公开演讲中
提出了 AI 产业的\textbf{五层架构}愿景，
这一框架成为理解 NVIDIA 战略布局的关键透镜。

\textbf{第一层：芯片与系统（Chips \& Systems）}。
NVIDIA 的根基：
从 GPU 芯片（H100/H200/B200/B300）
到完整的 AI 系统（DGX、HGX）、
到网络互联（NVLink、NVSwitch、Quantum InfiniBand）。
黄仁勋反复强调的一个观点是：
AI 计算不是「把 GPU 插到服务器上」这么简单，
芯片、内存、网络必须作为一个\textbf{系统}来设计和优化。
DGX SuperPOD 就是这种系统思维的产物：
它不是 GPU 的简单堆叠，
而是一个经过端到端优化的 AI 工厂。
截至 2026 年初，B200 的 FP8 算力是 H100 的 2.5 倍，
而 GB300 NVL72（72 块 GPU 的液冷系统）
则代表了 NVIDIA 在系统级集成上的最新能力。

\textbf{第二层：CUDA 与软件平台}。
CUDA 生态是 NVIDIA 最深厚的护城河，
超过 500 万开发者使用 CUDA，
几乎所有 AI 框架（PyTorch、TensorFlow、JAX）
都深度依赖 CUDA 内核。
这种软件锁定效应意味着
即使竞争对手的硬件在某些指标上追上 NVIDIA——
开发者的迁移成本仍然是巨大的障碍。
cuDNN、TensorRT、NCCL 等库
构成了一个完整的加速计算软件栈，
其成熟度远超任何竞品。

\textbf{第三层：AI 基础设施服务}。
DGX Cloud 和与云服务商的合作
使 NVIDIA 从芯片供应商扩展为基础设施服务商。
黄仁勋的愿景是：
NVIDIA 不仅卖 GPU，
还帮助客户运营「AI 工厂」：
提供从硬件到软件到优化服务的全套方案。

\textbf{第四层：AI Foundation Models}。
NVIDIA 通过 NeMo 框架、
NIM（NVIDIA Inference Microservices）
和 NVIDIA AI Enterprise 软件平台，
参与到模型训练和部署的全流程中。
这一层的战略意图是：
确保无论谁的模型最终胜出，
训练和部署都离不开 NVIDIA 的软件栈。

\textbf{第五层：行业应用}。
从自动驾驶（DRIVE）到医疗（Clara）
到机器人（Isaac）到数字孪生（Omniverse），
NVIDIA 正在向垂直行业渗透。
黄仁勋认为 AI 的最终价值
不在于模型本身，
而在于\textbf{改变每一个行业的生产流程}——
NVIDIA 希望在这个价值链的每一层都分一杯羹。

\begin{whybox}[NVIDIA 为什么（至今）无可替代？]
截至 2026 年，尽管有多个竞争者
（Google TPU、AMD MI300X、
Intel Gaudi 3、Broadcom 定制芯片、华为 Ascend），
NVIDIA 仍然占据 AI 训练市场 80\%+ 的份额。
原因不仅是硬件性能的领先
（B200 在多数基准上确实最快），
更重要的是\textbf{软件生态的锁定效应}：
CUDA 的开发者数量、
已有代码库的兼容性需求、
以及 NCCL 在分布式训练中的成熟度，
共同构成了一个极难打破的护城河。
竞争对手的策略因此分化为两条路线：
\textbf{兼容路线}（AMD 的 ROCm 试图兼容 CUDA 代码）
和\textbf{差异化路线}（Google TPU + JAX 建立独立生态）。
到目前为止，两条路线都只取得了有限的进展。
\end{whybox}

\section{训练成本的演变：2024--2026}
\label{sec:training-cost}

训练 frontier LLM 的成本正在快速上升，
但效率的提升也同样惊人。
理解这个双重趋势对于评估 AI 行业的可持续性至关重要。

\subsection{成本的量级}

截至 2026 年初，已知的训练成本估算如下：

\begin{table}[H]
\centering\small
\renewcommand{\arraystretch}{1.3}
\caption{Frontier 模型训练成本估算（2024--2026）}\label{tab:training-costs}
\begin{tabularx}{\linewidth}{l l l >{\raggedright\arraybackslash}X}
\toprule
\rowcolor{figLightGray}
\textbf{模型} & \textbf{训练时间} & \textbf{估算成本} & \textbf{备注} \\
\midrule
GPT-4 (2023) & $\sim$3 个月 & \$78M--100M+ &
  早期估算；含 A100 集群的 GPU 租赁成本 \\
Gemini Ultra (2023) & 数月 & \$191M（估算） &
  使用 Google 自有 TPU v4 集群 \\
LLaMA 3 405B (2024) & $\sim$54 天 & \$60M--100M（估算） &
  16,384 块 H100，Meta 自有集群 \\
DeepSeek-V3 (2024) & $\sim$2 个月 & \$5.5M（GPU 租赁） &
  2,048 块 H800；不含前期研发成本 \\
GPT-5 / Orion (2025) & 数月 & \$300M--500M（传闻） &
  未经官方确认 \\
下一代 frontier (2026) & --- & \$500M--1B+ &
  行业预估；含多轮训练 \\
\bottomrule
\end{tabularx}
\end{table}

DeepSeek-V3 报告的 \$5.5M 成本
是 AI 行业的一个标志性数字——
它证明了通过工程创新，
训练成本可以比同等性能的竞品低一个数量级。
但需要注意，这个数字仅包含最终训练的 GPU 租赁成本，
不包含数据处理、前期实验、后训练和人员成本。
完整的端到端成本可能在 \$20M--30M 左右。

\subsection{成本的构成}

训练一个 frontier LLM 的总成本可以分解为：

\begin{enumerate}[noitemsep]
  \item \textbf{GPU 计算}（50--70\%）：
        这是最大的成本项。
        以 2026 年初的价格，
        NVIDIA H100 的云端租赁价格约为 \$2.00--3.50/GPU$\cdot$小时，
        新一代 B200 约为 \$4.00--6.00/GPU$\cdot$小时。
        一次使用 8,192 块 H100 训练 60 天的成本：
        $8192 \times 60 \times 24 \times 2.50 \approx \$29.5M$。
  \item \textbf{数据获取与处理}（5--15\%）：
        网页爬取基础设施、数据标注
        （高质量 SFT 数据的标注成本约 \$15--50/对话）、
        数据清洗的计算成本。
  \item \textbf{人员}（15--25\%）：
        AI 研究员的年薪在硅谷通常为 \$300K--800K
        （含股权），顶尖人才超过 \$1M。
        一个 30 人的核心训练团队的年人员成本可达 \$15M--20M。
  \item \textbf{基础设施}（5--10\%）：
        网络设备、存储系统、电力和冷却。
        一个大规模 GPU 集群的年度电力成本可达数百万美元。
  \item \textbf{失败的实验}（不确定）：
        这是最难量化但可能占比很高的成本。
        训练到一半发现数据质量问题、
        超参数配置不佳导致 loss 不收敛，
        这些失败的尝试消耗的计算资源有时比最终成功的训练还多。
\end{enumerate}

\begin{corebox}[2026 年的成本趋势]
两个相反的力量正在塑造训练成本的走向：
\textbf{向上的力量}——模型规模和训练数据量的持续增长，
以及推理模型的 RL 训练需要更多的计算；
\textbf{向下的力量}——更高效的架构
（MoE 的稀疏激活将 FLOPs 降低 5--10 倍）、
更好的硬件（B200 的 FP8 性能是 H100 的 2.5 倍，单位算力价格持续下降）、
更聪明的训练方法（如 DeepSeek 的各种工程优化）。
净效果是：训练同等质量模型的成本在下降，
但 frontier 的定义不断被推高，
所以绝对成本仍在上升。
行业中流传一个说法：
“Training costs are dropping 10x per year for the same quality,
but we keep raising the bar 10x per year.”
\end{corebox}

\subsection{云端 vs.\ 自建集群的经济学}
\label{sec:cloud-vs-onprem}

选择使用云服务商的 GPU 还是自建集群，
是每个 AI 团队面临的关键战略决策。

\textbf{云端（Cloud）的优势}：
无需前期资本支出（CapEx），
按使用量付费（OpEx），
弹性伸缩（训练完成后释放资源），
以及云服务商提供的托管服务
（网络配置、故障自动替换、存储管理）。
主要的云 GPU 提供商包括
AWS（P5 instances with H100/H200, Trn2 instances with Trainium2）、
Google Cloud（A3 with H100, Cloud TPU v5p/v6e）、
Microsoft Azure（ND H100/H200 instances）
和 CoreWeave（GPU 密集型专用云）。
2026 年初，spot instance 价格可以将 H100 成本降低至
\$1.20--1.80/GPU$\cdot$小时，
但 spot instances 可能被中断，
对于需要持续运行数周的预训练任务不太适用。

\textbf{自建集群（On-Premise）的优势}：
长期来看成本更低——
Deloitte 的研究表明，大规模持续使用场景下
自建可以实现云端 60--70\% 的成本。
一块 H100 SXM 的采购价约为 \$25,000--35,000，
如果使用 3 年，每 GPU$\cdot$小时的折旧成本仅约 \$1.00--1.30
（不含电力、冷却、运维人员成本）。
此外，自建集群可以完全控制网络拓扑、
存储架构和安全策略，对于处理高度敏感数据的组织至关重要。

\textbf{关键的交叉点}：
根据 2026 年的 TCO（Total Cost of Ownership）分析，
当 GPU 利用率持续超过 60--70\% 且使用时间超过 18--24 个月时，
自建集群的经济性通常优于云端。
这就是为什么 Meta、Google、Tesla（xAI）等大型玩家
都选择自建，他们的 GPU 几乎 24/7 满载运行。
而初创公司和学术团队通常选择云端，
他们的使用是间歇性的，且无法承担前期的资本投入。

\textbf{混合策略}正在成为主流：
许多中型 AI 公司采用 “base on-prem + burst to cloud” 的策略，
自建一个满足日常需求的基础集群，
在需要大规模训练时临时租用云 GPU 扩展容量。

\begin{table}[H]
\centering\small
\renewcommand{\arraystretch}{1.3}
\caption{云端 vs.\ 自建集群：关键对比（2026 年）}\label{tab:cloud-vs-onprem}
\begin{tabularx}{\linewidth}{l >{\raggedright\arraybackslash}X >{\raggedright\arraybackslash}X}
\toprule
\rowcolor{figLightGray}
\textbf{维度} & \textbf{云端} & \textbf{自建} \\
\midrule
前期投入 & 极低 & 数百万--数千万美元 \\
单 GPU 成本/小时 & \$2.00--6.00（按需） & \$1.00--1.80（含运维） \\
弹性 & 分钟级扩缩容 & 数月的采购周期 \\
维护 & 服务商负责 & 需要专业运维团队 \\
数据安全 & 受限于服务商政策 & 完全自主 \\
3 年 TCO（1000 GPU） & \$52M--105M & \$35M--55M \\
适合谁 & 初创、学术、间歇性使用 & 大型实验室、持续训练 \\
\bottomrule
\end{tabularx}
\end{table}

\section{AI 人才市场}
\label{sec:talent-market}

AI 人才市场在 2024--2026 年经历了
前所未有的结构性变化——
薪资水平飙升、地理集中度加剧、
人才争夺日趋白热化。

\subsection{薪资与竞争}

2026 年的 AI 人才薪资已经达到了
让其他科技行业侧目的水平：

\begin{itemize}[noitemsep]
  \item \textbf{顶尖研究员}（发表过有影响力论文的高级研究者）：
        总薪酬（含股权）\$1M--5M/年。
        少数明星研究者（如从 OpenAI/Google 出走创业的）
        获得的 founding package 可达 \$10M+。
  \item \textbf{高级训练工程师}
        （有 frontier 模型训练经验的资深工程师）：
        总薪酬 \$500K--1.5M/年。
        这一薪资水平使得 AI 训练工程师成为
        硅谷薪酬最高的工程职位之一，
        超过了大多数 VP/Director 级别的管理岗位。
  \item \textbf{中级 ML 工程师}
        （3--5 年经验，能独立完成模型微调和部署）：
        总薪酬 \$250K--500K/年。
  \item \textbf{应届博士}（来自顶尖 AI 实验室的博士毕业生）：
        起薪 \$200K--350K/年。
        在竞争最激烈的时候
        （如 2024--2025 年 OpenAI、Anthropic、Google DeepMind 争夺应届生），
        签约奖金（signing bonus）可达 \$100K--300K。
\end{itemize}

在中国，顶尖 AI 人才的薪资同样在快速上升，
虽然绝对数字低于硅谷：
\begin{itemize}[noitemsep]
  \item 顶尖研究员（如 DeepSeek、月之暗面的核心成员）：
        年薪 200--500 万人民币，外加股权期权。
  \item 高级工程师：年薪 100--250 万人民币。
  \item 应届博士（清华、北大 AI 方向）：
        起薪 60--120 万人民币。
\end{itemize}

\subsection{地理集中与人才流动}

全球 AI 人才呈现高度的\textbf{地理集中}：

\textbf{旧金山湾区}仍然是全球 AI 人才密度最高的地区。
OpenAI、Anthropic、Meta FAIR、
Google DeepMind 的主要团队都位于此。
一个流传的说法是
「全球 AI 研究 top 100 的人，有 60 个住在旧金山」——
虽然数字不精确，但方向是对的。
这种集中产生了强烈的\textbf{聚集效应}
（agglomeration effect）：
人才之间的非正式交流
（咖啡馆对话、AI 沙龙、跳槽后的知识扩散）
创造了一个在其他地方难以复制的创新生态。

\textbf{北京}是中国 AI 人才的核心聚集地。
清华大学、北京大学、中科院
提供了持续的高质量人才供给。
DeepSeek、月之暗面、智谱AI、零一万物
均总部位于北京。
杭州（阿里巴巴）、深圳（腾讯、华为）
和上海（商汤、MiniMax）
是其他重要的 AI 人才中心。

\textbf{伦敦}是欧洲 AI 的枢纽
（DeepMind 总部、Mistral 的部分团队），
但其人才规模和薪资水平
与旧金山和北京相比仍有显著差距。

\textbf{Brain drain}（人才外流）模式值得关注：
\begin{itemize}[noitemsep]
  \item \textbf{学术界 → 产业界}：
        几乎是单向流动。
        产业界的薪资优势（3--10 倍于学术薪资）
        使得越来越多的顶尖教授离开大学
        或转为兼职，加入 AI 公司。
        这对学术 AI 研究的长期质量
        构成了结构性威胁。
  \item \textbf{中国 → 美国}（以及反向流动）：
        过去十年，大量中国 AI 博士
        留在美国的顶尖实验室工作。
        但 2024--2025 年出现了
        一波\textbf{回流潮}——
        DeepSeek、月之暗面等公司的崛起
        和有竞争力的薪资
        吸引了部分海外华人研究者回国。
        同时，美国的签证限制和地缘政治紧张
        也在推动部分中国背景的研究者
        寻找替代选择。
  \item \textbf{大公司 → 创业公司}：
        从 Google、Meta 等大公司
        流向 AI 创业公司的人才浪潮
        在 2024--2025 年达到高峰。
        Mistral、Sakana AI、SSI 等公司
        的创始团队几乎全部来自大型 AI 实验室。
\end{itemize}

\section{算力治理与地缘政治}
\label{sec:compute-governance}

AI 算力已经成为21世纪最重要的战略资源之一。
围绕 GPU 的出口管制、芯片供应链和算力基础设施的博弈，
正在深刻影响全球 AI 竞争的格局。

\subsection{美国对华出口管制的演进}

美国对中国的 AI 芯片出口管制
经历了三个逐步收紧的阶段：

\textbf{第一轮：2022 年 10 月}。
商务部工业与安全局（BIS）首次对
高性能 AI 芯片实施出口管制，
禁止向中国出口超过
\textbf{A100 级别}的芯片
（具体阈值基于算力和互联带宽）。
NVIDIA 随后推出了 A800 和 H800，
降规版（de-rated）芯片，
在管制阈值以下但仍然具有强大的训练能力。
DeepSeek-V3 就是在 H800 上训练的。

\textbf{第二轮：2023 年 10 月}。
管制范围扩大，
阈值进一步收紧至\textbf{H800 级别}也被禁止，
并将管制扩展到第三国转口贸易。
NVIDIA 再次推出降规版，\textbf{H20}——
这是一块几乎被「阉割」到只剩推理能力的芯片
（训练性能大幅下降，但推理吞吐量仍然可观）。

\textbf{第三轮：2025 年 1 月}。
拜登政府在卸任前发布了最全面的出口管制规则，
不仅针对芯片本身，还扩展到
\textbf{AI 模型权重}和\textbf{云服务}——
禁止将超过一定算力规模的 AI 训练服务
通过云端提供给中国用户。
这轮管制还引入了
\textbf{国家分级许可制度}（tiered licensing）：
将世界各国分为三个层级，
对盟友国家（日本、韩国等）保持开放，
对中国等国家严格限制。

\subsection{NVIDIA H20 的处境与替代方案}

\textbf{H20} 成为管制下的一个独特产品。
它的存在本身就是管制政策的产物——
NVIDIA 为了保住中国市场的收入
而专门设计的合规芯片。

H20 的关键参数：
\begin{itemize}[noitemsep]
  \item FP8 算力约为 H100 的 $\sim$30\%。
  \item HBM3 内存 96GB
        （与 H100 SXM 的 80GB 相比反而更大，
        这是因为 H20 的核心计算能力被削弱后，
        NVIDIA 在内存上做了补偿，
        以保持推理场景的竞争力）。
  \item 内存带宽 4TB/s（与 H100 的 3.35TB/s 相比更高）。
  \item NVLink 带宽被大幅削减。
\end{itemize}

H20 的定位是\textbf{推理芯片}而非训练芯片。
对于大规模预训练，H20 的效率远低于 H100，
训练同等模型需要约 3--5 倍的 H20 芯片。
但对于推理部署（特别是长上下文场景，
内存容量和带宽是瓶颈），
H20 仍然是中国市场上性能最强的合法选择之一。

\subsection{华为 Ascend 910C：国产替代的进展}

\textbf{华为 Ascend 910C}是中国最先进的
国产 AI 训练芯片。

关于 910C 的公开信息有限
（华为对芯片规格的披露非常谨慎），
但行业分析指出：
\begin{itemize}[noitemsep]
  \item FP16 算力约为 H100 的 60--70\%。
  \item 软件生态（CANN、MindSpore）
        与 NVIDIA 的 CUDA 生态相比仍有显著差距。
        但 2025--2026 年，
        华为大幅增加了对 PyTorch 兼容层的投入，
        降低了迁移成本。
  \item 互联带宽和网络拓扑
        （华为自研的 HCCS 互联）
        是与 NVIDIA 差距最大的领域，
        大规模分布式训练的效率
        受互联带宽的制约尤为严重。
\end{itemize}

DeepSeek 据报道已经在 V4 的开发中
进行华为 Ascend 的适配优化，
这不仅是技术选择，也是地缘政治现实下的必要准备。
如果未来连 H20 级别的芯片也被禁止出口，
中国 AI 实验室将不得不完全依赖国产替代方案。

\subsection{管制对中国 AI 的实际影响}

出口管制的影响是\textbf{双刃剑}。

\textbf{直接的负面影响}：
中国实验室无法获得最新的 NVIDIA GPU，
在原始算力上与美国实验室的差距在扩大。
一个 H100 集群能做到的事情，
用 H800 可能需要 1.5 倍的芯片，
用 H20 可能需要 3--5 倍，
用 Ascend 910C 可能需要 2--3 倍并承受更多的工程摩擦。
这种「算力税」在绝对规模竞争中是真实的劣势。

\textbf{意外的正面影响}——\textbf{约束驱动的创新}——
正是因为算力受限，
中国团队被迫在算法和工程效率上做到极致。
DeepSeek 的 MLA 架构将 KV cache 压缩了 10 倍以上，
DualPipe 和 DeepEP 优化了通信效率，
\$5.5M 的训练成本证明了
「用更少的资源做更多的事」不仅可能，
而且可以产生范式级的创新。
这种效率创新的外溢效应
甚至影响了美国实验室的研究方向——
DeepSeek-R1 发布后，
多个美国团队开始研究如何复现其低成本训练方法。

\textbf{芯片走私与灰色市场}也是不可忽视的现实。
媒体报道和政府调查表明，
尽管有出口管制，
仍有大量高端 NVIDIA 芯片
通过第三国转口进入中国~\cite{harithas2024chipsmuggling}。
BIS 在 2025 年加强了对东南亚
（特别是新加坡、马来西亚）
转口贸易的审查，
但完全阻止芯片流通在实践中极为困难。

\begin{warnbox}[出口管制的长期博弈]
出口管制的根本悖论在于：
短期内它限制了中国的 AI 能力，
但长期来看它可能\textbf{加速}中国的
芯片自主可控进程。
如果没有管制，中国实验室会继续依赖 NVIDIA，
因为 CUDA 生态的锁定效应极其强大。
管制强制打破了这种依赖，
为华为等国产芯片创造了
一个原本不可能的市场窗口。
2025--2026 年的证据表明，
中国正在快速缩小国产芯片的差距，
不是因为突然的技术突破，
而是因为管制创造了
「不用不行」的需求驱动力。
\end{warnbox}

\subsection{全球算力竞赛的量级}

将不同玩家的算力规模放在一起比较，
有助于理解竞争格局的全貌：

\begin{table}[H]
\centering\small
\renewcommand{\arraystretch}{1.3}
\caption{主要 AI 组织的 GPU 集群规模估算（2026 年初）}\label{tab:gpu-clusters}
\begin{tabularx}{\linewidth}{l l >{\raggedright\arraybackslash}X}
\toprule
\rowcolor{figLightGray}
\textbf{组织} & \textbf{GPU 规模（估算）} & \textbf{备注} \\
\midrule
Meta & $\sim$600K+ H100 等效 & 2026 年目标 1.3M GPU \\
Google & $\sim$500K+ TPU v5 等效 & 自研 TPU，不完全可比 \\
Microsoft/OpenAI & $\sim$300K+ H100 & Azure AI 集群 \\
xAI & $\sim$200K H100 & Memphis 超级集群 \\
Amazon & $\sim$150K+ 混合 & Trainium2 + H100 \\
DeepSeek & $\sim$10K--20K H800 & 受出口管制限制 \\
Anthropic & 公开信息有限 & 通过 AWS/GCP 租用 \\
Mistral & 数千块 GPU & 主要租用云端 \\
\bottomrule
\end{tabularx}
\end{table}

这张表揭示了一个关键事实：
在原始算力上，中国实验室与美国顶尖玩家
存在\textbf{一到两个数量级}的差距。
Meta 的 130 万 GPU 目标
是 DeepSeek 整个集群的 \textbf{60--100 倍}。
这种差距使得 DeepSeek 以 \$5.5M 训练出 V3
变得更加令人印象深刻，
它不得不在算法和工程效率上
做到远超竞争对手的极致水平
来弥补硬件的绝对劣势。

但算力不是一切。
历史上多次证明，
\textbf{算法创新可以抵消算力劣势}——
Transformer 本身就是一个例子
（它比 RNN 更高效地利用了并行计算资源）。
MoE 架构、MLA 注意力、
更好的数据配比和训练策略，
都在事实上放大了每单位算力的价值。
2025--2026 年的竞争格局表明，
\textbf{算力 $\times$ 效率 = 有效能力}——
两个因素同等重要。

\subsection{能源与环境挑战}

大规模 AI 训练的能源消耗
正在成为一个不可忽视的政策和商业问题。

一个 10,000 GPU 的 H100 集群
在满载运行时的功耗约为 \textbf{7--10 MW}——
相当于一个小型城镇的用电量。
加上冷却系统（传统风冷集群的 PUE 约 1.3--1.5，
液冷集群可降至 1.05--1.1），
一个大规模训练中心的年度电力成本
可达 \textbf{\$10M--30M}。

2025--2026 年，电力供应已经成为
AI 数据中心选址的\textbf{首要约束}。
北弗吉尼亚（全球最大的数据中心聚集地）
已经出现电网容量不足的问题，
新的 AI 数据中心不得不等待 2--3 年
才能获得足够的电力接入。
这催生了新的选址模式，
一些公司开始在
水电资源丰富的地区（如冰岛、挪威、加拿大魁北克）
或核电站附近建设数据中心，
以获得廉价且低碳的电力供应。

\textbf{液冷技术}的普及是对抗能耗增长的关键杠杆。
NVIDIA GB200 NVL72 系统已经\textbf{要求}液冷，
传统风冷无法满足其散热需求。
这标志着 AI 基础设施
从「服务器机房」向「工业冷却系统」的范式转变。
液冷不仅降低了 PUE（从 1.3+ 降至 1.1 以下），
还允许更高密度的 GPU 部署
（每个机架的 GPU 数量从 8 块增加到 72 块），
从而降低了数据中心的总体建设成本。

\section{开源训练框架生态}

不同的团队可以选择不同的训练框架，
各有侧重。

\subsection{预训练框架}

\begin{table}[H]
\centering\small
\renewcommand{\arraystretch}{1.3}
\caption{主流 LLM 预训练框架对比}\label{tab:training-frameworks}
\begin{tabularx}{\linewidth}{l l l >{\raggedright\arraybackslash}X}
\toprule
\rowcolor{figLightGray}
\textbf{框架} & \textbf{维护方} & \textbf{规模} & \textbf{核心优势} \\
\midrule
Megatron-LM & NVIDIA & 产业级 &
  最成熟的 3D 并行实现（TP + PP + DP），
  大多数主要实验室的训练起点 \\
DeepSpeed   & Microsoft & 产业级 &
  ZeRO 优化器~\cite{rajbhandari2020zero}，
  MoE 支持，与 HuggingFace 集成良好 \\
PyTorch FSDP & Meta (PyTorch) & 中大规模 &
  原生 PyTorch API，使用门槛低，
  但在超大规模上不如 Megatron 优化彻底 \\
TorchTitan  & Meta (PyTorch) & 中大规模 &
  LLaMA 训练的参考实现，
  基于 DTensor 分布式原语 \\
Nanotron    & HuggingFace & 研究级 &
  代码简洁可读，适合理解训练流程，
  教学和研究用途 \\
LitGPT      & Lightning AI & 中小规模 &
  易用的训练框架，
  支持多种模型架构，上手快 \\
\bottomrule
\end{tabularx}
\end{table}

\textbf{Megatron-LM} 是大规模预训练的事实标准。
它实现了完整的 3D parallelism（tensor、pipeline、data parallelism），
以及 sequence parallelism 和 expert parallelism 等扩展。
几乎所有训练过 100B+ 模型的团队都在 Megatron 的基础上构建。
它的缺点是代码复杂度高、学习曲线陡峭，
但对于需要在数千块 GPU 上稳定运行数月的训练任务，
没有更可靠的选择。

\textbf{DeepSpeed} 的核心贡献是 ZeRO 优化器，
它通过分片 optimizer states、gradients 和 parameters
大幅降低了每块 GPU 的内存需求。
DeepSpeed-Chat 提供了完整的 RLHF 训练流程支持。
在微调场景下，DeepSpeed 的 ZeRO-Offload
甚至允许在单块 GPU 上微调数十亿参数的模型
（通过将部分状态 offload 到 CPU 内存）。

\textbf{PyTorch FSDP}（Fully Sharded Data Parallel）
是 PyTorch 原生的分布式训练方案，
概念上类似于 DeepSpeed ZeRO Stage 3。
它的优势是与 PyTorch 生态无缝集成，
对于已有 PyTorch 代码库的团队迁移成本最低。

\subsection{后训练/微调框架}

对于后训练（SFT/RLHF/DPO），
生态系统更加多样化：

\begin{itemize}[noitemsep]
  \item \textbf{TRL}（HuggingFace）：
        集成 SFT、DPO、PPO、GRPO 训练，
        与 Transformers 和 PEFT 无缝对接。
        社区最活跃，文档最完善。
  \item \textbf{OpenRLHF}：
        专注于 RLHF/GRPO 训练，
        支持 Ray + vLLM 的大规模分布式 RL。
        在需要将训练规模扩展到数百块 GPU 时，
        OpenRLHF 的分布式 RL 架构更有优势。
  \item \textbf{LLaMA-Factory}：
        一键式微调工具，支持 100+ 模型，
        提供 Web UI 和命令行两种方式。
        特别适合快速实验和原型验证。
  \item \textbf{ms-swift}（阿里巴巴 ModelScope）：
        覆盖 SFT、RL、推理、部署的全流程工具链。
        与中国的模型生态（Qwen、通义千问等）集成最好。
  \item \textbf{Axolotl}：
        微调专用工具，支持多种方法
        （LoRA~\cite{hu2021lora}、QLoRA~\cite{dettmers2023qlora}、
        full fine-tuning、FSDP）。
        配置文件驱动，灵活度高。
  \item \textbf{Unsloth}：
        专注于极致的微调速度和内存优化。
        通过手写 Triton kernel 和自定义 autograd，
        在相同硬件上可以比标准实现快 2--5 倍。
\end{itemize}

\section{开源社区生态：从工具到运动}
\label{sec:opensource-ecosystem}

开源 LLM 不仅是一种技术选择，
更是一场改变 AI 领域权力结构的运动。
2025--2026 年，开源社区的成熟度和影响力达到了新的高度。

\subsection{HuggingFace：ML 领域的 GitHub}

HuggingFace 已经成为 ML 领域的核心基础设施，
类似于 GitHub 之于软件开发。
截至 2026 年初，其平台托管了超过 100 万个模型和
30 万个数据集，月活跃用户超过 500 万。

其生态系统包括：

\begin{itemize}[noitemsep]
  \item \textbf{Model Hub}：托管数十万个预训练模型，
        支持一行代码下载和使用。
        模型卡（Model Card）标准化了模型文档，
        包括训练数据、评估结果、已知限制和许可证信息。
  \item \textbf{Datasets}：标准化的数据集托管和加载，
        覆盖数千个 NLP/CV 数据集。
        FineWeb（15T tokens 的高质量网页数据）
        和 FineWeb-Edu（教育质量过滤的子集）
        是 HuggingFace 对开源训练数据的重要贡献。
  \item \textbf{Transformers 库}：
        统一的 API 支持几乎所有主流模型架构，
        是学术界和工业界最广泛使用的 NLP 库。
        新模型发布后通常在数天内就会被集成。
  \item \textbf{TRL + PEFT}：
        后训练和参数高效微调的标准工具。
        TRL 在 2025 年加入了 GRPO 支持，
        使得复现 DeepSeek-R1 的训练流程成为可能。
  \item \textbf{SmolLM 系列}：
        HuggingFace 自研的小模型家族
        （135M、360M、1.7B 参数），
        展示了数据质量在小模型训练中的关键作用。
        SmolLM2~\cite{allal2025smollm2} 通过精心策划的训练数据
        （Cosmopedia 合成教科书、FineWeb-Edu、Stack 代码数据）
        在同等规模下达到了领先性能。
  \item \textbf{Open LLM Leaderboard}：
        自动化的模型评估排行榜，
        是社区衡量模型质量的重要参考。
        2025 年升级为 v2 版本，
        引入了更难的评估基准（GPQA、MUSR、BBH）
        以应对模型能力的快速提升。
  \item \textbf{Spaces}：
        免费托管模型 demo 的平台，
        研究者可以零成本展示自己的工作。
\end{itemize}

HuggingFace 的核心贡献是\textbf{降低了 ML 的准入门槛}。
在 HuggingFace 之前，使用一个新模型需要：
阅读论文 → 找到作者的 GitHub repo → 配置环境 → 下载权重
→ 编写推理代码。
在 HuggingFace 之后，这简化为三行 Python 代码。

\subsection{本地推理生态}

llama.cpp 和 GGUF 格式的出现
创造了一个完整的本地推理生态系统。
llama.cpp 使用纯 C/C++ 实现，
支持 CPU、Apple Silicon（Metal）、CUDA、Vulkan 等多种后端，
使得在没有 NVIDIA GPU 的消费级硬件上运行 LLM 成为可能。

\textbf{Ollama} 和 \textbf{LM Studio}
进一步将本地推理的门槛降到了非技术用户可及的水平，
一键安装、一键下载模型、开箱即用的聊天界面。
这使得「人人都能在自己的电脑上运行 AI」
从技术理想变成了现实。

社区围绕量化模型建立了活跃的生态：
模型发布后通常在数小时内就会出现
GGUF 格式的各种量化版本
（Q2\_K、Q4\_K\_M、Q5\_K\_S、Q8\_0 等），
用户可以根据自己的硬件条件选择合适的量化级别。

\subsection{社区模型：百花齐放}

开源社区的模型产出已经形成了一个完整的层次结构：

\textbf{第一梯队：实验室级开源}——
DeepSeek-V3/R1（MIT 许可证）、
LLaMA 4~\cite{dubey2024llama3}（Llama 许可证）、
Qwen 3~\cite{qwen2025qwen3}（Apache 2.0）、
Gemma 3（Google，开放权重）。
这些模型的质量已经接近甚至达到闭源模型的水平。

\textbf{第二梯队：完全开源}——
OLMo 2（AI2，数据+代码+权重全部公开）、
SmolLM2（HuggingFace，完整训练流程可复现~\cite{allal2025smollm2}）、
MAP-Neo（M-A-P 社区，7B 全开源 bilingual 模型）。
这些项目的核心价值不在于模型性能的绝对领先，
而在于\textbf{可复现性}——
研究者可以在完全相同的条件下重复实验。

\textbf{第三梯队：社区微调}——
基于第一梯队的 base model，
社区产出了大量针对特定场景的微调模型。
仅在 HuggingFace 上，
基于 LLaMA 的微调模型就超过 10 万个。
这种长尾效应使得开源生态的实际影响力
远超任何单一模型。

\begin{whybox}[开源 LLM 的「飞轮效应」]
开源 LLM 生态形成了一个自我强化的飞轮：
更好的模型 → 更多的用户 → 更多的微调和改进
→ 更活跃的社区 → 更好的工具和基础设施
→ 更低的使用门槛 → 更多的用户 → \ldots
HuggingFace、llama.cpp、Ollama 等工具的成功
证明了这个飞轮的力量。
一个关键观察是：
开源模型的性能与闭源模型的差距
从 2023 年的「代际差距」
缩小到 2025--2026 年的「几个月的差距」，
DeepSeek-R1 在某些数学推理任务上
甚至超越了同期的闭源竞品。
\end{whybox}

\subsection{Chatbot Arena：社区驱动的评估}

LMSYS 的 Chatbot Arena 是开源社区
对 LLM 评估最重要的贡献之一。
它使用类似于国际象棋 Elo 评分的系统，
让真实用户对不同模型的回答进行盲评。
由于评估者不知道哪个回答来自哪个模型，
且样本量足够大（数百万次投票），
Chatbot Arena 的排名被认为是
最能反映真实用户偏好的评估方法，
比任何单一 benchmark 都更可信。

LMArena 从 UC Berkeley 的学术项目
发展为独立运营的评估平台公司，
2025 年底获得超过 1 亿美元融资，
这本身就说明了
评估作为 AI 基础设施的战略价值。
在一个「模型即商品」的时代——
\textbf{评估的可信度就是权力}——
谁定义了「好」的标准，
谁就在很大程度上影响了
市场的竞争规则和资源的分配方向。

\subsection{开源许可证的战略考量}

开源 LLM 的许可证选择
反映了各组织截然不同的战略意图：

\textbf{MIT / Apache 2.0}
（DeepSeek-V3/R1、Qwen 3）：
最宽松的许可证，允许任何使用
包括商业用途，几乎没有限制。
这种选择的战略意图是
\textbf{最大化生态系统采纳率}：
让尽可能多的人使用你的模型，
建立技术标准和社区影响力。

\textbf{Llama License}
（Meta LLaMA 系列）：
非标准的定制许可证，
允许大多数商业使用
但对月活超过 7 亿的产品施加限制。
这个阈值精确地瞄准了
Meta 的直接竞争对手
（Google、Amazon、ByteDance），
同时对其他所有人保持开放。
这是一种\textbf{精算式的开源}——
看似慷慨，实则精心设计了
排除条款来保护竞争优势。

\textbf{完全开源}
（AI2 OLMo、HuggingFace SmolLM）：
不仅开放权重，还公开训练数据、
训练代码和所有实验细节。
这种做法的价值不在于模型性能的领先，
而在于\textbf{科学可复现性}：
使得独立研究者可以完全复现训练过程，
验证声明的结果，
并在此基础上进行有意义的改进。

\section{从论文到产品：被低估的工程挑战}

在学术界和媒体报道中，
人们往往关注算法创新和基准测试分数，
而忽略了将 LLM 从实验室推向产品所需的大量工程工作。

\subsection{训练基础设施的可靠性}

在 2048 块 GPU 上训练两个月，
意味着 $2048 \times 60 \times 24 = 295$ 万 GPU$\cdot$小时。
如果每块 GPU 每天有 0.1\% 的故障率，
那么平均每小时会有 $\sim$2 块 GPU 出现故障。
训练系统必须能够自动检测故障、
隔离故障节点、从最近的检查点恢复训练。
DeepSeek-V3~\cite{deepseekv3} 的训练报告强调他们实现了
“zero irrecoverable loss spikes”，
这不是运气，而是精心设计的容错机制的结果。

\subsection{调试分布式训练}

当 loss 在第 50,000 步突然发散，而你用了 4096 块 GPU，
怎么找到 bug？
这是分布式训练中最令人恐惧的场景之一。
可能的原因：某块 GPU 的 ECC 内存错误导致梯度异常；
某个数据 shard 包含异常样本（如超长文本或乱码）；
通信超时导致部分梯度丢失；
甚至可能是浮点数非确定性，
相同的 checkpoint 在不同的 GPU 配置下恢复训练，
可能产生不同的 loss 曲线。

实践中的调试策略包括：
详尽的日志系统（记录每一步的 loss、gradient norm、
learning rate、每块 GPU 的计算时间）；
“canary” 训练（在小集群上用相同配方并行运行，
比较 loss 曲线是否一致）；
以及最重要的——\textbf{可重现性}，
或者更准确地说，对不可重现性的接受和管理。

\textbf{故障隔离}（fault isolation）是大规模调试的核心方法论。
当 loss spike 出现时，
工程师需要快速缩小问题范围：
首先检查是否是\textbf{硬件故障}——
通过比较每块 GPU 的 gradient norm 分布
可以识别产生异常梯度的节点
（正常节点的 gradient norm 应当高度一致）；
其次排查\textbf{数据问题}，
回溯 spike 发生时正在处理的数据 shard，
检查是否包含异常样本（极长文本、编码错误、
重复数据导致的梯度方向偏移）；
最后考虑\textbf{数值稳定性}，
混合精度训练中的 FP16 溢出、
allreduce 操作中的精度损失等。
2025--2026 年的实践经验表明，
约 60\% 的训练中断源自硬件故障，
30\% 来自数据问题，
仅 10\% 是算法或超参数配置错误。

\textbf{分布式 profiling 工具}的成熟度
直接决定了调试的效率。
PyTorch 的内置 profiler、
NVIDIA Nsight Systems 和 Ray 的分布式追踪
可以精确定位通信瓶颈和计算/通信的重叠效率。
但在 4096+ GPU 规模下，
即使是 profiling 数据本身的收集和分析
也是一个不小的工程挑战，
单次 profiling 可能产生数 TB 的追踪数据。

\subsection{数据管道的效率}

2048 块 GPU 的训练每秒消耗数百万 token，
数据加载不能成为瓶颈。
这需要高效的数据格式（如 memory-mapped 二进制文件）、
多级缓存（SSD → CPU 内存 → GPU 内存）
和异步预取（在 GPU 计算时提前准备下一批数据）。

数据管道的设计远比表面看起来复杂。
\textbf{预处理阶段}需要对原始数据（通常是 PB 级的网页爬取数据）
进行去重、质量过滤、个人信息脱敏和 tokenization。
这些步骤本身就需要数千 CPU 核心并行运行数天，
且必须保证确定性，
因为训练的可重现性要求相同的数据、以相同的顺序、
产生完全相同的 token 序列。
Meta 在 2025 年底公开的 Data PreProcessing Service（DPP）
展示了 exabyte 级训练集群的数据管道架构：
通过流式 tokenization 和分级缓存
消除了数据加载的 GPU 空闲等待（data stall）。

\textbf{数据混合的动态调度}是另一个关键工程挑战。
frontier 模型的训练数据通常由多个来源混合而成，
网页文本、书籍、代码、数学、对话等，
每个来源的采样权重不同。
更先进的做法是在训练过程中\textbf{动态调整}这些权重：
根据持续评估的能力曲线反馈，
自动增加表现不足的能力维度对应的数据权重。
这要求数据管道支持在线的权重变更而无需中断训练：
一个看似简单但工程上极具挑战性的需求，
因为它涉及分布式系统中数千个数据加载 worker 的同步协调。

\textbf{存储层的选择}直接影响训练吞吐量和成本。
PCIe Gen5 NVMe SSD（2025 年起成为训练集群标配，
单盘顺序读取超过 14 GB/s）使得本地缓存成为可行方案，
而高性能并行文件系统（如 WEKA、Lustre、GPFS）
则解决了跨节点的数据共享问题。
WEKApod 方案已实现 8 个存储节点为 768 块 H100 GPU
提供 720 GB/s 的持续读取带宽，
这种吞吐量足以确保即使在最大规模的训练中，
数据管道也不会成为瓶颈。

\subsection{人的因素}

最后一个常被忽略的挑战是\textbf{人}。
一次 frontier model 的训练可能消耗数千万美元的计算资源。
如果训练在第六周因为一个 bug 失败而需要重启，
这不仅是经济损失，也是巨大的心理压力。
负责训练的工程师需要 7$\times$24 小时 on-call，
凌晨 3 点收到 loss spike 的告警是常态。
燃尽（burnout）是 LLM 训练团队面临的真实挑战。

2026 年初，OpenAI 一位资深工程师公开辞职时
坦言「持续的 on-call 压力和对训练失败的焦虑
严重影响了他的心理健康」，
这一事件引发了 AI 行业对工程师福祉的广泛讨论。
一项 2026 年的调查显示，
约 70\% 的 SRE（Site Reliability Engineer）
报告了严重的 alert fatigue（告警疲劳），
而 AI 训练团队的 on-call 强度远超传统 SRE，
因为一次训练中断的经济损失
可能相当于普通生产事故的数百倍。

\textbf{组织层面的缓解措施}正在逐步建立。
成熟的训练团队采用\textbf{轮班制 on-call}
（至少 3--4 人的轮值队列，确保每人每月不超过一周 primary on-call），
并投资于\textbf{自动化恢复系统}：
能够自动检测故障、回退 checkpoint、隔离问题节点的系统
可以将需要人工介入的事件减少 80\% 以上。
DeepSeek 的「zero irrecoverable loss spikes」
不仅是技术成就，也是对团队心理压力的巨大缓解。
一些领先的团队还引入了\textbf{心理健康支持}：
包括训练项目结束后的强制休假期、
匿名压力报告机制、以及专业心理咨询资源。
这些举措反映了一个日益清晰的认知：
在 frontier AI 训练中，
\textbf{团队的可持续性与模型的可持续性同样重要}。

这些工程挑战解释了为什么
只有少数组织能够成功训练 frontier LLM，
即使算法和数据都是公开的，
将它们在大规模上可靠地运行仍然是一项极其困难的工程任务。
更关键的是，支撑这些工程系统的\textbf{人}，
他们的经验、判断力和心理韧性，
才是真正不可复制的稀缺资源。

\subsection{模型发布的安全流程}

将训练好的模型推向用户
需要经过一套严格的\textbf{安全发布流程}，
这个流程在 2025--2026 年逐步标准化。

\textbf{分阶段发布}（Staged Release）是主流实践：
\begin{enumerate}[noitemsep]
  \item \textbf{内部测试}（Alpha）：
        仅限内部员工使用。
        重点测试功能完整性和基本安全性。
        通常持续 2--4 周。
  \item \textbf{有限预览}（Beta）：
        邀请外部的信任合作伙伴和安全研究者使用。
        红队在此阶段进行深入的安全评估。
        收集早期用户反馈。持续 2--6 周。
  \item \textbf{公开预览}（Preview/Early Access）：
        向更广泛的用户开放，
        但保留撤回或修改的权利。
        密切监控使用模式和安全事件。
  \item \textbf{正式发布}（General Availability）：
        向所有用户开放。
        附带完整的系统卡（System Card），
        详细记录模型的能力、限制、
        已知安全风险和推荐使用方式。
\end{enumerate}

\textbf{系统卡}（System Card / Model Card）
已成为负责任发布的标准组成部分。
一份好的系统卡应包含：
模型的训练数据概况（来源、规模、时间范围）、
在主要基准上的表现、
已知的安全风险（如偏见、幻觉、越狱脆弱性）、
不推荐的使用场景
（如医疗诊断、法律建议等高风险场景）、
以及安全评估的结果摘要。
GPT-4、Claude 3 和 Gemini 的系统卡
都遵循了这一框架，
虽然各家在信息披露的完整度上有差异。

\textbf{开源模型的安全责任}是一个更复杂的问题。
当 DeepSeek 或 Meta 将模型权重完全公开时，
他们无法控制下游用户如何使用这些权重：
包括移除安全防护（un-aligning）、
微调用于有害用途等。
这种「一旦发布就失去控制」的特性
使得开源模型的安全发布
需要更多地依赖\textbf{事前评估}
（发布前确认模型不具备关键的危险能力）
而非\textbf{事后控制}
（通过 API 限制和使用政策约束用户行为）。

\section{2025--2026 年的人才与格局变迁}

2025 年下半年至 2026 年初，
AI 领域的人才流动和组织重组以前所未有的烈度展开。
两个标志性事件勾勒出了中国 AI 竞争的新轴线。

\subsection{姚顺雨加入腾讯}

2025 年 12 月 17 日，27--28 岁的\textbf{姚顺雨}（Shunyu Yao）
被任命为腾讯首席 AI 科学家（Chief AI Scientist），
直接向总裁刘炽平汇报。
姚顺雨此前是 OpenAI 研究员，
以 Tree of Thoughts、ReAct 等工作闻名，
是 Agent 领域最具影响力的年轻研究者之一。

腾讯为此次任命做了大规模的组织重构，
同步新设三个一级部门：
AI 基础设施部、AI 数据部和数据算力平台部。
姚顺雨直接统管 AI 基础设施部和大模型部，
获得了罕见的跨部门资源调度权。
在此之前，腾讯已于 12 月 8 日发布了
\textbf{混元 2.0}（Hunyuan 2.0）：
一个 406B/32B MoE 模型，支持 256K 上下文，
为姚顺雨的到来做了技术铺垫。

姚顺雨的研究哲学值得关注：
他多次强调「\textbf{方法的复杂度应匹配任务的难度}」
（method complexity should match task difficulty），
这是一种反对 benchmark 军备竞赛、
追求方法优雅性的研究品味。
他规划的下一代模型将以 Agent 能力为核心，
参数规模约 30B，
刻意选择了效率优先而非规模优先的路径。
朗桥证券的分析报告指出：
「姚顺雨的蜜月期可能只有 6 个月」，
腾讯需要在 2026 年下半年看到明确的技术突破，
否则这场高调引进的战略价值将面临质疑。

\subsection{梁文锋与 DeepSeek 团队}

如果说 2025 年 1 月的 DeepSeek-R1 是「中国的 Sputnik 时刻」，
那么 2025 年 9 月 R1 登上 \textit{Nature} 封面
则是对这一成就的正式学术认证。
\textbf{梁文锋}（Wenfeng Liang）作为通讯作者，
以不到 200 人的团队完成了 frontier 级别的研究突破。

更令人惊叹的是 DeepSeek 团队在 6 个月内的论文产出：
R1（\textit{Nature}，2025 年 9 月）、
mHC（arXiv:2512.24880，2025 年 12 月 31 日，梁文锋共同作者）、
Engram（arXiv，2026 年 1 月 12 日，梁文锋共同作者），
三篇论文覆盖了推理训练、架构创新和记忆系统三个方向，
展示了小团队在研究广度上的惊人能力。
DeepSeek APP 的月活用户已突破 1 亿。

V4 的开发据报道聚焦于三个方向：
原生多模态、LTM（Long-Term Memory）突破、
以及国产芯片（华为 Ascend）适配优化。
但梁文锋面临的压力也是空前的：
在 R1 引发全球关注之后，
DeepSeek 必须再次交出
「世界最好的开源模型」这一答卷，
而这一次竞争对手们已经做好了充分准备。

\subsection{竞争格局的两极化}

梁文锋与姚顺雨成为 2026 年中国 AI 的双焦点，
但两人面临的挑战截然不同。

梁文锋是「所有人都在等待他交出下一个作品的明星创业者」——
DeepSeek 必须证明 R1 的成功不是偶然，
V4 需要在技术上超越已经高度内卷的竞争对手，
同时应对来自出口管制和芯片限制的地缘政治压力。

姚顺雨是「从硅谷空降、被委以改革重任的 95 后明星科学家」——
他需要在腾讯这个以产品见长但在 AI 基础研究上相对滞后的组织中，
建立全新的技术文化，
在政治资本耗尽之前拿出让人信服的成果。

但竞争格局远不止这两个焦点人物。
\textbf{阿里巴巴 Qwen 团队}在 2026 年初经历了一次痛苦的人才流失。
2026 年 3 月，Qwen 大模型后训练负责人\textbf{郁博文}
离职加入字节跳动 Seed 团队，
负责视觉模型和多模态交互方向的后训练工作。
在此之前，通义实验室大模型技术负责人林俊旸也已离职。
这些核心人才的流失并非孤立事件，
阿里巴巴在 2025 年底进行了大规模组织重构，
将 AI 资源向「Alibaba Token Hub」和
企业级 AI 平台「悟空」集中，
而这种战略重心的转移导致了基础研究团队的不稳定。
尽管如此，Qwen3 系列（包括原生多模态的 Qwen3-Omni）
仍然是中国开源模型中最具竞争力的选手之一。

\textbf{字节跳动 Seed 团队}正在通过积极的人才吸纳
构建多模态 AI 的全面能力。
从 Qwen 团队挖来的后训练专家、
从海外回流的研究人员、
以及内部从豆包产品线积累的实战经验，
使 Seed 团队在 2026 年成为中国 AI 人才密度最高的团队之一。
行业分析指出，顶尖人才的选择
越来越取决于两个因素：
\textbf{技术愿景的一致性}
和\textbf{算力资源的可获得性}，
而字节跳动在这两方面都具有强大的吸引力。

\textbf{国际维度}使竞争格局更加复杂。
在美国，OpenAI、Anthropic、Google DeepMind 和 Meta
继续在 frontier 模型上保持领先地位，
但中国团队通过效率创新
（如 DeepSeek 的低成本训练范式）
和开源策略不断缩小差距。
出口管制限制了中国团队获取最新 NVIDIA GPU 的能力，
但也倒逼了国产芯片生态的加速发展
和训练效率的极致优化，
这种「约束驱动的创新」
正在产生意料之外的技术溢出效应。
欧洲则以 Mistral AI 为代表，
试图在美中两极之间建立第三条路径，
但在算力和人才规模上仍面临显著差距。

两种模式——
小团队极致效率 vs.\ 巨头组织重构，
谁能在 2026 年下半年率先交出答卷，
将深刻影响中国乃至全球 AI 产业的走向。
但更深层的趋势是，
竞争正在从\textbf{模型能力}转向\textbf{生态构建}——
谁能将模型嵌入最多的应用场景、
建立最强的数据飞轮、
吸引最多的开发者和用户，
谁就能在下一阶段的竞争中占据优势。
纯粹的基准分数竞赛正在让位于
更复杂的系统性竞争。

\subsection{全球 AI 产品生态的演变}

在模型层之外，AI 产品层的竞争同样激烈，
2025--2026 年的产品生态经历了深刻的重构。

\textbf{AI 编程工具的并购整合}。
\textbf{Cognition AI} 在 2025 年 7 月收购了 Windsurf，
时间恰在 Google 以 24 亿美元
反向收购 Windsurf CEO Varun Mohan 和核心团队之后。
这一连串事件揭示了 AI 编程工具市场的整合逻辑：
自主 Agent（Devin）需要 IDE 的交互体验，
IDE 需要 Agent 的自主执行能力，
两者的融合是不可避免的。
Cognition 在 2025 年 9 月完成 4 亿美元融资，
估值达到 102 亿美元：
一个成立仅两年的 Agent 公司
获得了与成熟 SaaS 企业可比的估值，
体现了资本市场对 Agent 赛道的极端看好。

\textbf{GitHub Copilot} 从代码补全工具
进化为完整的 Agent 平台。
2025 年 4 月 Agent mode 面向所有 VS Code 用户推出，
2026 年 1 月引入 Plan mode（Shift+Tab 切换），
2 月新增模型选择器（每个任务选择最适合的模型），
3 月在 JetBrains IDE 中全面 GA。
Copilot 的独特优势是与 GitHub 生态的深度集成，
Agent 可以直接创建 PR、触发 Actions 工作流、
管理 issue，将代码编写和项目管理统一在一个界面中。

\textbf{Cursor} 以 AI-first IDE 的定位
在开发者中获得了狂热追随。
其 Agent Mode 可以自主分析代码库、
进行多文件编辑、运行终端命令，
2026 年 3 月引入的 Subagents 功能
允许多个专业化 Agent 在主 Agent 控制下并行运行。
Cursor 的商业模式（订阅制 IDE + API 调用费用）
正在挑战传统 IDE 的免费+插件模式。

\textbf{AI 搜索产品的崛起}。
\textbf{Perplexity} 以 AI 原生搜索引擎的定位
在 2025--2026 年获得了快速增长。
OpenAI 的 SearchGPT 和 Google 的 AI Overviews
代表了传统搜索巨头的 AI 化转型。
这三条路线的竞争
可能从根本上改变信息检索的范式：
从“十个蓝色链接”到“一个结构化答案”。

\textbf{中国 AI 产品生态}。
豆包（字节跳动）、Kimi（月之暗面）、DeepSeek Chat、
通义千问（阿里巴巴）在中国市场形成了激烈竞争。
DeepSeek APP 的月活用户已突破 1 亿，
Kimi 以超长上下文和 Agent 能力
在专业用户群中建立了差异化优势。
中国 AI 产品市场的特点是
\textbf{价格战驱动的快速渗透}，
几乎所有产品都提供免费或极低价的基础服务，
通过使用量和企业服务变现。

\textbf{AI 基础设施公司}。
Together AI、Fireworks AI、Modal、Replicate 等
AI 基础设施公司在 2025--2026 年快速增长。
它们提供从模型部署到推理优化的全栈服务，
填补了“拥有开源模型权重”
和“能以生产级可靠性服务用户”之间的鸿沟。
OpenRouter 作为 LLM 路由器，
在多模型比价和流量分发中扮演了关键角色，
其统计数据已成为行业分析的重要参考。

\textbf{从“模型为王”到“应用为王”的转变}。
2024 年，AI 产品的核心竞争力是接入最好的模型；
2026 年，当多个 frontier 模型
以接近零的成本可用时，
竞争力转向了上层，
\textbf{工作流设计}（如何将 AI 能力嵌入用户的工作流程）、
\textbf{数据飞轮}（如何利用用户交互数据改进体验）、
\textbf{领域知识}（如何针对特定行业定制解决方案）。
这个转变正在重新定义 AI 公司的护城河，
模型能力不再是壁垒，
用户体验和领域深度才是。

\subsection{OpenAI 创始团队的衰减曲线（2026 Q2）}

2026 年 4 月 17 日，
OpenAI 一日内宣布三位核心高管同时离职：
前 CPO、OpenAI for Science 负责人 \textbf{Kevin Weil}，
Sora 团队负责人 \textbf{Bill Peebles}，
以及企业应用 CTO \textbf{Srinivas Narayanan}~\cite{openai_three_exits_2026}。
两年内，11 位联合创始人中仅 Sam Altman 和 Greg Brockman 留任；
仅 2025 年就有至少 12 位高管离开。
4 月初，CPO Fidji Simo 因神经免疫疾病休病假，
产品线由 Brockman 临时代管；
CMO Kate Rouch 因癌症康复离任，
COO Brad Lightcap 调入“特殊项目”。
Sora 因日均 100 万美元的算力亏损在 3 月底关停，
4 月 26 日彻底下线；
OpenAI for Science 解散并入 Codex 团队。

人才三向分流的格局在 2026 年春季已经清晰：
（1）追求安全研究与稳定文化的人流向 \textbf{Anthropic}；
（2）追求顶级现金包裹的人流向 \textbf{Meta MSL}；
（3）希望保留股权上行空间的人则集中流向 \textbf{Anysphere}（Cursor）、
\textbf{Cognition}、\textbf{Thinking Machines Lab} 等创业公司。
研究员留存率的差异印证了这一分化：
Anthropic 据报“压倒性领先”，
DeepMind 78\%，OpenAI 仅 67\%~\cite{brin_claude_taskforce_2026}。
DeepMind 已启用 6--12 个月带薪竞业条款延缓人才外流。

\subsection{Meta MSL 与超级薪酬包反向博弈}

Meta Superintelligence Labs（MSL）在 2025 年下半年至 2026 年初
以前所未有的现金强度重塑人才市场。
2026 年 4 月 21 日，MSL 一次性挖走了
Mira Murati 创立的 \textbf{Thinking Machines Lab} 五位创始成员，
其中 Andrew Tulloch 的包裹据报“潜在数十亿美元”%
~\cite{meta_msl_thinking_machines_2026}。
同月早些时候，MSL 还从 OpenAI 挖走 4 名研究员
（包括 Shengjia Zhao、Jiahui Yu）。

但反向跳槽案例同样耀眼。
\textbf{Ruoming Pang} 在 2025 年中
被 MSL 以 2 亿美元以上包裹从 Apple 挖走，
仅 7 个月后于 2026 年 2 月 25 日跳槽 OpenAI~\cite{pang_meta_to_openai_2026}。
OpenAI 普通研究员现金加股权约 150 万美元/年，
Pang 包裹未披露但行业普遍预期“远超此数”。
这是 MSL 巨额签约金\textbf{反向博弈}的标志性案例：
当 OpenAI 需要留人时，
即使是 Meta 的“天花板包裹”也可以被超越。

MSL 的早期产出仍未跟上其投入。
4 月 8 日发布的首个原生多模态推理模型 \textbf{Muse Spark}
保持闭源~\cite{wang_msl_muse_spark_2026}；
内部代号 \texttt{Avocado} 的第二个模型据报仍落后 Gemini 3 Pro。
Yann LeCun 在 2025 年 11 月离职后曾公开称 Alexandr Wang “年轻、缺乏经验”。
Meta 2025 年全年营收 \$201B（同比 +22\%），
2026 年 capex 指引 \$115--135B，
其中包括与 Nebius 合资 \$27B 的 GW 级数据中心~\cite{meta_nebius_jv_2026}。

\subsection{中国大厂 AI 人才环流图（2026 Q1/Q2）}

中国 AI 顶层人才在 2026 年春季完成了一次大规模再分配。

\textbf{字节跳动 Seed} 已成为中国版 MSL：
DeepSeek-Coder 与 R1 主力作者、95 后 L8 工程师 \textbf{郭达雅}
4 月以“近亿元年薪”加入 Seed；
此前 Qwen 后训练负责人 \textbf{余博文} 3 月加盟 Seed；
Qwen Code 前负责人惠斌 1 月赴 Meta；
原通义实验室总负责人林俊旸离职~\cite{deepseek_talent_outflow_2026}。
Qwen 团队的后训练职责已交由前 DeepMind 高级研究员周昊接管。

\textbf{DeepSeek 人才外流}是 2026 上半年中国 AI 圈最被低估的故事：
LLM 早期成员王炳轩转入腾讯，
DeepSeek-OCR 一作魏昊然离职，
Janus-Pro 一作阮冲加入 DeepRoute.ai，
R1 训练成员罗福莉加入小米。
梁文锋本人于 2025 年 6 月底
向 DeepMind VP Wu Yonghui 发出顶级 offer 但未能成功。
员工期权缺位是核心结构性问题：
High-Flyer 量化基金 2025 年收益 56.6\%，
但 DeepSeek 员工无法参与上行，
长期内卷下大批核心成员选择套现离场~\cite{deepseek_external_funding_2026}。

\textbf{Sergey Brin 亲自督战编码追赶}是另一个标志。
2026 年 4 月，Brin 与 CTO Koray Kavukcuoglu 共同督导 Google DeepMind 的
“Claude 追赶组”，由前 DeepMind 预训练主管
\textbf{Sebastian Borgeaud} 领导，
明确目标是“用 Gemini 编码反超 Claude Code、
自动化 Google 自身工作流”~\cite{brin_claude_taskforce_2026}。
Hassabis 与 Pichai“每天通话”，
Google 2026 年整体 AI 投入约 \$185B/年。
在 Brain $\times$ DeepMind 合并三年后，
Hassabis 已成为 6,000 人组织实质上的 CEO 候选，
Jeff Dean 转任首席科学家、操作权重下降。

\subsection{Anthropic：估值阶梯式狂飙与三巨头算力—投资循环}

Anthropic 的估值与现金流在三个月内
经历了 AI 史上最陡峭的一段曲线。
2026 年 2 月 Series G 募集 \$30B，post-money \$380B；
4 月中旬 VC 抢入价跃至 \$800B，
Amodei 兄妹尚未接受~\cite{anthropic_800b_valuation_2026}。
ARR 演进：2024 年底 \$1B → 2025 年 12 月 \$9B
→ 2026 年 3 月 \$30B，
是企业软件史上从 \$1B 到 \$30B 用时最短的纪录。
Claude Code 主导开发者工具市场，
企业客户占总收入的 80\%（OpenAI 仅 40\%）。

围绕 Anthropic 的算力—投资循环已经形成：
（1）2026 年 4 月 24 日 \textbf{Google} 承诺最高 \$40B
（首期 \$10B 基于 \$350--380B 估值），
配套 5 年 5GW 的 Google Cloud 算力承诺~\cite{google_anthropic_40b_2026}。
（2）2026 年 4 月 20 日 \textbf{亚马逊}追加最高 \$25B
（即时 \$5B + 里程碑 \$20B），
Anthropic 承诺 10 年向 AWS 采购 \$100B 算力，
锁定 Trainium2/3/4 至 2027 年~\cite{aws_anthropic_25b_2026}。
（3）微软 2025 年 11 月已投入 \$5B，
绑定 \$30B Azure 采购合约。
亚马逊累计承诺达 \$33B，
持股估算 16--20\%；
Google 持股 14\%，累计投入 \$3B，
账面已实现 \$10.7B 未实现收益。
这种“三巨头通过对 Anthropic 的投资换取算力消纳”的模式，
是 2026 年 AI 资本结构最重要的创新——
hyperscaler 不再是中立的算力供应商，
它们与前沿实验室的命运深度绑定。

\subsection{OpenAI \$852B 私募与 IPO 2027}

2026 年 3 月 31 日，OpenAI 完成 \$122B 一级私募，
post-money 估值 \$852B，
Amazon、Nvidia、SoftBank、Microsoft 共同领投，
首度向散户开放约 \$3B 额度~\cite{openai_852b_round_2026}。
CFO Sarah Friar 锁定 2027 年 IPO 时间窗，
\$4.7B 主承销由 JPM、Citi、GS、MS、WF 五行包销~\cite{openai_ipo_2027}。
OpenAI 与 AWS 续签 \$100B/8 年合作，
承诺至少消耗 2GW AWS 算力。
ARR 已突破 \$25B（39 个月达成）；
按此估算估值/ARR 倍数约 34$\times$，
显著高于历史上任何一家私有 SaaS 的同期水平。

\subsection{xAI 并入 SpaceX：从独立公司到子公司形态}

2026 年 2 月 2 日，
\textbf{SpaceX 全股票收购 xAI}，
xAI 转为 SpaceX 全资子公司~\cite{xai_spacex_merger_2026}。
交易估值 SpaceX \$1T、xAI \$250B，
合体 \$1.25T。
2026 年 1 月 xAI 还独立完成 Series E \$20B
（\$230B post-money），
现已并入母公司财务结构。
Grok 5 仍在训练，
但 SpaceX 工程师内部更倾向使用 Anthropic Claude，
Musk 推动以 \$10B 合作或 \$60B 收购方式吸收 \textbf{Cursor}，
并已从 Cursor 挖走两名高级员工。
SpaceX 同时从 Starlink 调入 Michael Nicolls 整顿 xAI 工程组织。
xAI 章节的叙事自此发生根本变化：
组织上并入 SpaceX，
资本上脱离独立公司形态，
其前途与 Musk 整体商业帝国的命运绑定。

\subsection{DeepSeek 首次外部融资}

2026 年 4 月 17 日，
The Information 首次披露 DeepSeek 启动外部融资，
目标至少 \$300M、估值 \$10B 起%
~\cite{deepseek_external_funding_2026}。
4 月 21--22 日谈判升级到\textbf{腾讯、阿里}参与，
估值预期超 \$20B（对标 MiniMax \$40B）。
腾讯有意拿下 20\% 股权，
但梁文锋拒绝大幅稀释。
背景压力直接：
V4 多次延期 15 个月，
训练 Ascend 出现稳定性问题被迫回退 NVIDIA H800；
员工期权缺失导致核心团队大批出走。
首次外部融资意味着 DeepSeek 从“独立的 High-Flyer 内部研究项目”
转向\textbf{标准 AI 创业公司治理结构}，
这是其文化转型的关键一步。

\subsection{应用层独角兽估值跃迁}

2026 年 4 月，应用层公司估值进入新区间：

\begin{itemize}[noitemsep]
  \item \textbf{Cursor (Anysphere)} 在谈 \$2B+，
        pre-money \$50B（2025 年 11 月 Series D 还是 \$29.3B），
        ARR 2026 年 2 月已达 \$2B，
        Fortune 1000 企业中 70\% 是其客户~\cite{cursor_50b_round_2026}。
  \item \textbf{Cognition} 在谈数亿美元，估值 \$25B
        （2025 年 9 月 \$10.2B 翻倍多），
        Devin 2.0 价格从 \$500/月 降到 \$20/月（$-96\%$），
        3 月以约 \$250M 收购 Windsurf~\cite{cognition_25b_round_2026}。
  \item \textbf{Mistral} 2026 年 3 月 30 日完成 \$830M 首次债务融资
        （7 行欧洲银行团），
        购置 13,800 张 GB300 建巴黎 44MW 数据中心~\cite{mistral_debt_2026}。
  \item \textbf{Harvey} 3 月 25 日由 GIC + Sequoia 领投 \$200M、
        估值 \$11B（2025 年 12 月 \$8B），ARR \$190M~\cite{harvey_11b_round_2026}。
  \item \textbf{Perplexity} 2026 年 4 月 ARR 达 \$500M。
  \item \textbf{Sierra}（Bret Taylor）2025 年 9 月已达 \$10B。
\end{itemize}

“2024 年的独角兽溢价”
正在演变为“2026 年的十角兽（Decacorn）常态”。
ARR/估值倍数普遍位于 20--40$\times$ 区间，
高于传统 SaaS 历史上任何一个 vintage。

\subsection{H200“中间档”与 75K 上限：分级出口管制}

2026 年 1 月 15 日，
美国 BIS 终稿规则将 NVIDIA H200 与 AMD MI325X
从“presumption of denial”转为“case-by-case”审批，
TPP $<$ 21,000、DRAM 带宽 $<$ 6,500 GB/s 进入审批门槛%
~\cite{bis_h200_tiered_2026}。
3 月 17 日 NVIDIA 重启 H200 对华生产，
首批 82,000 GPU 准备发货，
每枚需经 BIS 实验室复核加 25\% 关税；
3 月报道特朗普考虑对中国买家施加 75K 颗/家上限~\cite{nvidia_h200_china_2026}。
北京一度让海关阻挡 H200，
后政策反转。
Blackwell 仍被全面禁止；
\textbf{B30A}（H20 性能 12 倍的中国特供）的批准与否
已成为 2026 年中美 AI 谈判的核心议题。

管制叙事自此修订为
\textbf{“分级关税 + 单笔许可 + 收入分成”}的三轨混合体——
不再是简单的“准卖/不准卖”，
而是按性能档位、买家规模、地缘风险三维度
动态调整的复杂博弈。

\subsection{华为 Ascend 出货量与中国混合栈}

2026 年华为 Ascend 出货预计 60 万枚 \textbf{910C}
（较 2025 年翻倍），
加上其他 Ascend 型号合计 160 万 die
（2025 年仅 100 万）~\cite{huawei_ascend_910c_volume_2026}。
910C 由 SMIC 增强 7nm DUV 多重曝光产出，
BF16 约 670 TFLOPS（约为 B200 的 1/3），
CANN 软件栈逐步成熟。
DeepSeek R2 因 Ascend 训练不稳被迫回退 NVIDIA H800，
中国“\textbf{训练用 N、推理用 Ascend}”
已成既定的混合算力栈结构。
这种分工与西方“训练推理同构在 NVIDIA 上”的范式形成鲜明对比，
也是中国 AI 基础设施在出口管制夹缝中的一次
\textbf{结构性适应}。

\subsection{Stargate 的两个支点：美国 9GW 与 UAE 5GW}

Stargate 项目在 2026 年春季显形为两个独立支点：

（1）\textbf{美国本土 9GW}。
Epoch AI 2026 年 4 月 20 日报告：
7 个站点已敲定（德克萨斯 3 个、新墨西哥、威斯康星、密歇根、俄亥俄），
合计超过 9GW，
约等于纽约市峰值用电，
可支撑约 2,000 万张 H100 等价算力~\cite{stargate_9gw_2026}。
Abilene 旗舰已 0.6GW 运行，
Oracle 6 月起交付 GB200 机柜，
OpenAI 已用其训练 + 推理。
4.5GW Oracle 增量加 5 个新址使总容量逼近 9--10GW、
投入超 \$450B。
SoftBank 2025 年 12 月 28 日收购 DigitalBridge（\$108B 数字基础设施），
但 The Information 2026 年 2 月披露 SoftBank 实际未派人参建，
OpenAI 改与 Oracle、SoftBank 分别直签。

（2）\textbf{阿布扎比 5GW}。
G42 + OpenAI + Oracle + Nvidia + Cisco + SoftBank 联合，
2026 年 4 月 21 日 G42 子公司 Khazna 公布：
首期 200MW 末期机电安装，
5,000+ 工人在岗，
目标 2026 Q3 投运%
~\cite{stargate_uae_5gw_2026}。
完整园区 5GW、19.2 km$^2$。
同期美国商务部批准向 G42 与沙特 HUMAIN 出口 70,000 张 GB300。
G42 CEO 在达沃斯设定 KPI：
“2026 年生产 10 亿个 AI agent，
每天 12 小时工作消耗约 1GW 电力。”
韩国宣布加入 5GW UAE 项目。
Stargate 自此从“美国本土项目”
扩展为\textbf{“跨大西洋—海湾—远东三极的算力联盟”}。

\subsection{大模型同价竞赛与单集群 GW 时代}

DeepSeek 触发的价格连锁反应在 2026 年春季全面显形。
阿里 Qwen API 从 \$1.10/M tokens
砍至 \$0.07（$-93\%$）跟进 DeepSeek~\cite{qwen_price_war_2026}；
Qwen 系列已超越 Llama 成为 HuggingFace 下载量第一的开源家族。
中国开源模型在 HuggingFace 全球开源用量份额
从 2024 年底 1.2\% 跃至 2026 年初约 30\%。
字节豆包（Doubao）2.0 周活 1.55 亿。
Qwen 3.5-397B 在 GPQA Diamond 取得 88.4\%，
反超 GPT-5.2、Claude Opus 4.5、Gemini 3 Pro。
应用层同样剧烈：
Devin 2.0 从 \$500/月 降到 \$20/月（$-96\%$）。

算力侧则进入\textbf{单集群 GW 时代}。
xAI \textbf{Colossus I/II} 已就位，
2026 Q1 后扩容；
Musk 称 Grok 5 训练用集群 $>$20 万 GPU 等效。
\textbf{Meta} 与 Nebius \$27B 合资建 GW 级 AI 数据中心~\cite{meta_nebius_jv_2026}，
并锁定 6.6 GW 核电（含 Three Mile Island 旁合作）。
\textbf{Microsoft} 通过 \$1.6B/20 年 PPA + DOE \$1B 贷款
重启 TMI Crane（约 835MW）；
\textbf{Google} $\times$ Kairos Power 500MW SMR；
\textbf{Amazon} \$500M 投资 X-energy（2039 年 5GW SMR）+
Susquehanna 1,920MW PPA 至 2042 年。
FERC 2025 年 12 月裁定 AI 设施可直挂核电厂、
绕过电网排队。
2026 年全球数据中心年耗电首次突破 1,000 TWh，
北弗吉尼亚高压变压器交期从 30 周拉到 120+ 周。
hyperscaler 已签核电承诺合计 $\geq$ 11.8 GW%
~\cite{datacenter_nuclear_ppas_2026}——
\textbf{电力本身}成为 frontier AI 的稀缺生产要素。

\section{组织文化与管理哲学}
\label{sec:org-culture}

观察不同 AI 实验室的成功模式，
\textbf{组织文化}的差异
对技术产出的影响
可能不亚于算力和人才的差异。

\subsection{「研究优先」vs.「产品优先」}

AI 实验室的核心文化选择
往往决定了它们的发展轨迹：

\textbf{研究优先}的组织
（如早期的 OpenAI、DeepMind、DeepSeek）
以发表突破性论文和推进科学边界为首要目标。
这种文化吸引了最优秀的研究者
（他们追求学术声誉和技术挑战），
但可能导致产品化缓慢，
因为「让论文更优雅」和「让用户更满意」
有时是矛盾的目标。
DeepMind 在被 Google 收购后
经历了从纯研究到产品交付的痛苦转型，
多年间发表了大量顶会论文
但在消费产品上的影响力远不如 OpenAI。

\textbf{产品优先}的组织
（如字节跳动 Seed、腾讯混元、AWS）
以用户价值和商业指标为首要目标。
这种文化擅长快速迭代和规模化交付，
但可能在基础研究上缺乏深度：
当你的 OKR 是用户增长而非论文发表时，
投入六个月研究一个可能失败的新架构
在激励机制上是不合理的。

\textbf{最成功的组织}往往是
两种文化的\textbf{精心混合}。
Anthropic 在安全研究（学术导向）
和 Claude 产品（市场导向）之间
建立了一种共生关系，
安全研究的突破成为产品差异化的核心卖点，
产品收入反哺安全研究的持续投入。
OpenAI 从 2023 年开始在内部
明确区分了 Research 和 Applied 两个轨道，
允许研究者不受产品时间线的约束，
同时确保 Applied 团队有足够的资源
将研究成果快速产品化。

\subsection{决策速度与组织层级}

在快速变化的 AI 领域，
\textbf{决策速度}是一种关键的组织能力。

DeepSeek 的成功部分归因于其\textbf{极度扁平}的组织结构。
据报道，梁文锋直接参与到技术细节层面，
他不仅是管理者，也是研究的积极参与者
（以共同作者身份出现在多篇论文中）。
这种「技术创始人深度参与」的模式
使得从实验发现到决策执行的路径极短，
但也高度依赖于创始人的个人带宽，
随着 DeepSeek 的规模和影响力增长，
这种模式的可扩展性将面临考验。

相比之下，Google DeepMind 的决策链条更长：
研究者的实验发现需要经过组长、
部门负责人、副总裁层层审批，
最终才能影响产品决策。
这种流程在避免错误决策方面更有保障，
但也可能导致「机会窗口已关闭
但审批还在进行中」的情况。

\textbf{War room 模式}
是大型训练项目中常见的临时组织形式：
在关键的训练期间（通常持续数周到数月），
将所有相关团队（研究、Infra、数据、评估）
集中在同一物理空间或永久在线的通信频道中，
取消日常的层级汇报结构，
实行扁平化的即时协作。
LLaMA 3 和 DeepSeek-V3 的训练
都采用了这种模式——
它在短期内极大地提升了决策速度和执行效率，
但代价是参与者的高强度工作压力。

\subsection{知识管理与技术传承}

训练 frontier 模型的知识
很大程度上是\textbf{默会知识}（tacit knowledge）——
它存在于工程师的直觉和经验中，
难以完整地写入文档。

「当 loss 在第 30,000 步出现微小的波动时，
应该调整学习率还是检查数据？」
「当 gradient norm 比预期高 5\% 时，
是否预示着未来的训练不稳定？」
这些判断无法通过阅读论文获得，
只能通过\textbf{亲身参与}多次大规模训练来积累。
这就是为什么有 frontier 训练经验的工程师
如此稀缺和昂贵，
全球可能只有数百人
拥有完整的、端到端的 100B+ 模型训练经验。

\textbf{知识传承}的挑战在于：
当一个关键工程师离职时，
他带走的不仅是技能，
还有大量无法文档化的\textbf{上下文知识}，
「为什么我们用这个学习率而不是那个？」
「上次出现类似的 loss spike 时我们做了什么？」
成熟的团队通过以下方式缓解这一风险：
\textbf{详尽的训练日志}
（记录每次训练的所有决策及其原因）、
\textbf{pair training}
（每次训练至少两人全程参与，互为备份）、
\textbf{内部技术分享}
（定期的 post-mortem 和经验交流会议）。
但即便如此，
关键人员的流失仍然是
AI 实验室面临的最大组织风险之一。

\begin{corebox}[训练 Frontier LLM 的组织要素总结]
成功训练一个 frontier LLM 需要以下要素的\textbf{同时满足}：
\begin{enumerate}[noitemsep]
  \item \textbf{算力}：数千到数万块顶级 GPU，稳定运行数月。
  \item \textbf{数据}：PB 级的高质量多样化训练数据。
  \item \textbf{人才}：数十到数百名具备
        frontier 训练经验的研究员和工程师。
  \item \textbf{资金}：数千万到数亿美元的持续投入。
  \item \textbf{文化}：研究深度与工程效率的平衡，
        快速决策与质量保障的平衡。
  \item \textbf{时间}：至少 6--12 个月的持续迭代。
\end{enumerate}
能够同时满足所有这些条件的组织，
在全球范围内不超过 20 个。
这就是为什么 frontier AI 训练
本质上是一项\textbf{组织竞争}，
而非纯粹的技术竞争。
\end{corebox}


%% ============================================================
%% BibTeX entries referenced in this chapter
%% (to be added to the main .bib file)
%%
%% @article{deepseekv3,
%%   title   = {{DeepSeek-V3} Technical Report},
%%   author  = {DeepSeek-AI},
%%   journal = {arXiv preprint arXiv:2412.19437},
%%   year    = {2024},
%% }
%%
%% @article{deepseekr1,
%%   title   = {{DeepSeek-R1}: Incentivizing Reasoning Capability in {LLMs} via Reinforcement Learning},
%%   author  = {DeepSeek-AI},
%%   journal = {arXiv preprint arXiv:2501.12948},
%%   year    = {2025},
%% }
%%
%% @article{penedo2024fineweb,
%%   title   = {{FineWeb}: decanting the web for the finest text data at scale},
%%   author  = {Penedo, Guilherme and Kydl{\'\i}{\v{c}}ek, Hynek and {Ben Allal}, Loubna and others},
%%   journal = {arXiv preprint arXiv:2406.17557},
%%   year    = {2024},
%% }
%%
%% @article{kaplan2020scaling,
%%   title   = {Scaling Laws for Neural Language Models},
%%   author  = {Kaplan, Jared and McCandlish, Sam and Henighan, Tom and others},
%%   journal = {arXiv preprint arXiv:2001.08361},
%%   year    = {2020},
%% }
%%
%% @article{hoffmann2022chinchilla,
%%   title   = {Training Compute-Optimal Large Language Models},
%%   author  = {Hoffmann, Jordan and Borgeaud, Sebastian and Mensch, Arthur and others},
%%   journal = {arXiv preprint arXiv:2203.15556},
%%   year    = {2022},
%% }
%%
%% @article{dubey2024llama3,
%%   title   = {The {Llama 3} Herd of Models},
%%   author  = {Dubey, Abhimanyu and Jauhri, Abhinav and Pandey, Abhinav and others},
%%   journal = {arXiv preprint arXiv:2407.21783},
%%   year    = {2024},
%% }
%%
%% @article{bai2022constitutional,
%%   title   = {Constitutional {AI}: Harmlessness from {AI} Feedback},
%%   author  = {Bai, Yuntao and Kadavath, Saurav and Kundu, Sandipan and others},
%%   journal = {arXiv preprint arXiv:2212.08073},
%%   year    = {2022},
%% }
%%
%% @article{touvron2023llama,
%%   title   = {{LLaMA}: Open and Efficient Foundation Language Models},
%%   author  = {Touvron, Hugo and Lavril, Thibaut and Izacard, Gautier and others},
%%   journal = {arXiv preprint arXiv:2302.13971},
%%   year    = {2023},
%% }
%%
%% @article{rajbhandari2020zero,
%%   title   = {{ZeRO}: Memory Optimizations Toward Training Trillion Parameter Models},
%%   author  = {Rajbhandari, Samyam and Rasley, Jeff and Ruwase, Olatunji and He, Yuxiong},
%%   journal = {arXiv preprint arXiv:1910.02054},
%%   year    = {2019},
%% }
%%
%% @article{hu2021lora,
%%   title   = {{LoRA}: Low-Rank Adaptation of Large Language Models},
%%   author  = {Hu, Edward J and Shen, Yelong and Wallis, Phillip and others},
%%   journal = {arXiv preprint arXiv:2106.09685},
%%   year    = {2021},
%% }
%%
%% @article{dettmers2023qlora,
%%   title   = {{QLoRA}: Efficient Finetuning of Quantized Language Models},
%%   author  = {Dettmers, Tim and Pagnoni, Artidoro and Holtzman, Ari and Zettlemoyer, Luke},
%%   journal = {arXiv preprint arXiv:2305.14314},
%%   year    = {2023},
%% }
%%
%% @article{jiang2024mixtral,
%%   title   = {Mixtral of Experts},
%%   author  = {Jiang, Albert Q and Sablayrolles, Alexandre and Roux, Antoine and others},
%%   journal = {arXiv preprint arXiv:2401.04088},
%%   year    = {2024},
%% }
%%
%% @article{anil2023gemini,
%%   title   = {Gemini: A Family of Highly Capable Multimodal Models},
%%   author  = {Anil, Rohan and Dai, Andrew M and Firat, Orhan and others},
%%   journal = {arXiv preprint arXiv:2312.11805},
%%   year    = {2023},
%% }
%%
%% @article{qwen2025qwen3,
%%   title   = {{Qwen3} Technical Report},
%%   author  = {Qwen Team},
%%   journal = {arXiv preprint},
%%   year    = {2025},
%% }
%%
%% @article{allal2025smollm2,
%%   title   = {{SmolLM2}: When Smol Goes Big --- Data-Centric Training of a Small Language Model},
%%   author  = {{Ben Allal}, Loubna and Lozhkov, Anton and Bakouch, Elie and others},
%%   journal = {arXiv preprint arXiv:2502.02737},
%%   year    = {2025},
%% }
